%%%%%%%%%%%%%%%%%%%%%%%%%%%%%%%%%%%%%%%%%%%%%%%%%%%%%%%%%%%%
% 8.13 MATHEMATICAL STATISTICS
% The Composition of Advanced Mathematics
%%%%%%%%%%%%%%%%%%%%%%%%%%%%%%%%%%%%%%%%%%%%%%%%%%%%%%%%%%%%

\section{Mathematical Statistics}
\label{sec:mathematical-statistics}

\prerequisite{
Probability theory, random variables, expectation and variance,
probability distributions, joint distributions, transformations,
sampling distributions, estimation, confidence intervals, hypothesis
testing, likelihood, regression, multivariate statistics, and calculus.
}

\learningobjectives{
By the end of this section, the reader should be able to:
\begin{itemize}
    \item describe a statistical model mathematically;
    \item distinguish parameters, parameter spaces, statistics, and
          estimators;
    \item define identifiability;
    \item define point estimators and estimator properties;
    \item distinguish bias, variance, and mean squared error;
    \item define consistency;
    \item understand asymptotic normality;
    \item derive estimators using the method of moments;
    \item derive maximum likelihood estimators;
    \item compute score functions and Fisher information;
    \item state and use the Cramér--Rao lower bound;
    \item define sufficient statistics;
    \item apply the factorization theorem;
    \item define minimal sufficiency conceptually;
    \item understand ancillary statistics;
    \item understand complete statistics;
    \item state the Rao--Blackwell theorem;
    \item state the Lehmann--Scheffé theorem;
    \item understand exponential families;
    \item identify natural parameters and sufficient statistics in
          exponential-family models;
    \item distinguish exact and asymptotic inference;
    \item understand likelihood-ratio, Wald, and score procedures;
    \item understand the Neyman--Pearson lemma;
    \item define most powerful and uniformly most powerful tests;
    \item understand monotone likelihood ratios;
    \item understand pivotal quantities;
    \item construct exact confidence intervals from pivots;
    \item understand confidence-set inversion;
    \item define loss functions and statistical risk;
    \item understand decision rules;
    \item understand admissibility and minimax ideas conceptually;
    \item distinguish frequentist and Bayesian decision rules;
    \item understand asymptotic efficiency;
    \item understand the delta method;
    \item understand Slutsky's theorem in statistical applications;
    \item recognize order-statistic theory;
    \item understand likelihood asymptotics;
    \item connect mathematical statistics with Bayesian statistics,
          regression, statistical learning, and computational statistics.
\end{itemize}
}

%%%%%%%%%%%%%%%%%%%%%%%%%%%%%%%%%%%%%%%%%%%%%%%%%%%%%%%%%%%%
\subsection{What Is Mathematical Statistics?}
\label{subsec:what-is-mathematical-statistics}
%%%%%%%%%%%%%%%%%%%%%%%%%%%%%%%%%%%%%%%%%%%%%%%%%%%%%%%%%%%%

Mathematical statistics develops the theoretical foundations of
statistical inference using probability and analysis.

Probability begins with a model and studies the behavior of random
quantities.

Statistics begins with observed data and attempts to learn about the
model.

Symbolically,

\[
\boxed{
\text{Probability: model}
\longrightarrow
\text{data behavior}
}
\]

while

\[
\boxed{
\text{Statistics: data}
\longrightarrow
\text{model information}.
}
\]

%%%%%%%%%%%%%%%%%%%%%%%%%%%%%%%%%%%%%%%%%%%%%%%%%%%%%%%%%%%%
\subsection{Statistical Models}
\label{subsec:statistical-models}
%%%%%%%%%%%%%%%%%%%%%%%%%%%%%%%%%%%%%%%%%%%%%%%%%%%%%%%%%%%%

A statistical model is a collection of probability distributions indexed
by one or more unknown parameters.

A parametric model may be written

\[
\boxed{
\mathcal{P}
=
\{
P_\theta:
\theta\in\Theta
\}.
}
\]

Here,

\[
\Theta
\]

is the parameter space.

%%%%%%%%%%%%%%%%%%%%%%%%%%%%%%%%%%%%%%%%%%%%%%%%%%%%%%%%%%%%
\subsection{Parameter Space}
\label{subsec:parameter-space}
%%%%%%%%%%%%%%%%%%%%%%%%%%%%%%%%%%%%%%%%%%%%%%%%%%%%%%%%%%%%

\begin{definitionbox}{Parameter Space}
The parameter space \(\Theta\) is the set of all parameter values allowed
by the statistical model.
\end{definitionbox}

For a Bernoulli model,

\[
\boxed{
\Theta=[0,1].
}
\]

For a normal model with unknown mean and variance,

\[
\boxed{
\Theta
=
\{
(\mu,\sigma^2):
\mu\in\mathbb{R},
\ \sigma^2>0
\}.
}
\]

%%%%%%%%%%%%%%%%%%%%%%%%%%%%%%%%%%%%%%%%%%%%%%%%%%%%%%%%%%%%
\subsection{Parameter Versus Statistic}
\label{subsec:parameter-versus-statistic}
%%%%%%%%%%%%%%%%%%%%%%%%%%%%%%%%%%%%%%%%%%%%%%%%%%%%%%%%%%%%

A parameter is a fixed but unknown quantity describing the population.

A statistic is a function of the observed sample only.

For example,

\[
\boxed{
\mu
=
\text{population parameter},
}
\]

while

\[
\boxed{
\overline{X}
=
\text{sample statistic}.
}
\]

Before observing the data, the statistic is itself random.

%%%%%%%%%%%%%%%%%%%%%%%%%%%%%%%%%%%%%%%%%%%%%%%%%%%%%%%%%%%%
\subsection{Estimator and Estimate}
\label{subsec:estimator-and-estimate}
%%%%%%%%%%%%%%%%%%%%%%%%%%%%%%%%%%%%%%%%%%%%%%%%%%%%%%%%%%%%

\begin{definitionbox}{Estimator}
An estimator is a statistic used to estimate an unknown parameter.
\end{definitionbox}

For example,

\[
\boxed{
\widehat{\mu}
=
\overline{X}.
}
\]

The estimator is the random variable

\[
\overline{X}.
\]

The estimate is the numerical value obtained after observing the sample.

%%%%%%%%%%%%%%%%%%%%%%%%%%%%%%%%%%%%%%%%%%%%%%%%%%%%%%%%%%%%
\subsection{Identifiability}
\label{subsec:statistical-identifiability}
%%%%%%%%%%%%%%%%%%%%%%%%%%%%%%%%%%%%%%%%%%%%%%%%%%%%%%%%%%%%

A parameter is identifiable if distinct parameter values correspond to
distinct probability distributions.

Formally,

\[
\boxed{
P_{\theta_1}
=
P_{\theta_2}
\Longrightarrow
\theta_1=\theta_2.
}
\]

If two parameter values generate exactly the same distribution, the data
cannot distinguish them.

Thus identifiability is necessary for meaningful estimation.

%%%%%%%%%%%%%%%%%%%%%%%%%%%%%%%%%%%%%%%%%%%%%%%%%%%%%%%%%%%%
\subsection{Sampling Distribution}
\label{subsec:sampling-distribution-mathematical-statistics}
%%%%%%%%%%%%%%%%%%%%%%%%%%%%%%%%%%%%%%%%%%%%%%%%%%%%%%%%%%%%

If

\[
T=T(X_1,\ldots,X_n)
\]

is a statistic, its probability distribution under repeated sampling is
called its sampling distribution.

For example, if

\[
X_i
\overset{\text{i.i.d.}}{\sim}
N(\mu,\sigma^2),
\]

then

\[
\boxed{
\overline{X}
\sim
N
\left(
\mu,
\frac{\sigma^2}{n}
\right).
}
\]

Mathematical statistics studies these distributions to evaluate
estimators and tests.

%%%%%%%%%%%%%%%%%%%%%%%%%%%%%%%%%%%%%%%%%%%%%%%%%%%%%%%%%%%%
\subsection{Bias}
\label{subsec:estimator-bias}
%%%%%%%%%%%%%%%%%%%%%%%%%%%%%%%%%%%%%%%%%%%%%%%%%%%%%%%%%%%%

\begin{definitionbox}{Bias}
The bias of an estimator \(\widehat{\theta}\) is

\[
\boxed{
\operatorname{Bias}_\theta
(
\widehat{\theta}
)
=
\mathbb{E}_\theta
[
\widehat{\theta}
]
-
\theta.
}
\]
\end{definitionbox}

An estimator is unbiased if

\[
\boxed{
\mathbb{E}_\theta
[
\widehat{\theta}
]
=
\theta
}
\]

for every \(\theta\in\Theta\).

%%%%%%%%%%%%%%%%%%%%%%%%%%%%%%%%%%%%%%%%%%%%%%%%%%%%%%%%%%%%
\subsection{Worked Example: Bias of the Sample Mean}
\label{subsec:worked-bias-sample-mean}
%%%%%%%%%%%%%%%%%%%%%%%%%%%%%%%%%%%%%%%%%%%%%%%%%%%%%%%%%%%%

\begin{workedexamplebox}{Unbiasedness of the Sample Mean}

Suppose

\[
X_1,\ldots,X_n
\]

are i.i.d. with

\[
\mathbb{E}[X_i]=\mu.
\]

Then

\[
\mathbb{E}[\overline{X}]
=
\mathbb{E}
\left[
\frac1n
\sum_{i=1}^{n}
X_i
\right].
\]

By linearity,

\[
\mathbb{E}[\overline{X}]
=
\frac1n
\sum_{i=1}^{n}
\mu
=
\mu.
\]

Therefore,

\begin{answerbox}
\[
\boxed{
\operatorname{Bias}(\overline{X})=0.
}
\]
\end{answerbox}

\end{workedexamplebox}

%%%%%%%%%%%%%%%%%%%%%%%%%%%%%%%%%%%%%%%%%%%%%%%%%%%%%%%%%%%%
\subsection{Variance of an Estimator}
\label{subsec:variance-estimator}
%%%%%%%%%%%%%%%%%%%%%%%%%%%%%%%%%%%%%%%%%%%%%%%%%%%%%%%%%%%%

The variance

\[
\boxed{
\operatorname{Var}
(
\widehat{\theta}
)
}
\]

measures sampling variability of the estimator.

Among unbiased estimators of the same parameter, lower variance is
generally preferred.

However, a slightly biased estimator can have much lower total error than
an unbiased estimator.

%%%%%%%%%%%%%%%%%%%%%%%%%%%%%%%%%%%%%%%%%%%%%%%%%%%%%%%%%%%%
\subsection{Mean Squared Error}
\label{subsec:mean-squared-error}
%%%%%%%%%%%%%%%%%%%%%%%%%%%%%%%%%%%%%%%%%%%%%%%%%%%%%%%%%%%%

\begin{definitionbox}{Mean Squared Error}
The mean squared error of \(\widehat{\theta}\) is

\[
\boxed{
MSE_\theta
(
\widehat{\theta}
)
=
\mathbb{E}_\theta
[
(
\widehat{\theta}
-
\theta
)^2
].
}
\]
\end{definitionbox}

It decomposes as

\[
\boxed{
MSE
=
\operatorname{Var}
(
\widehat{\theta}
)
+
\operatorname{Bias}
(
\widehat{\theta}
)^2.
}
\]

%%%%%%%%%%%%%%%%%%%%%%%%%%%%%%%%%%%%%%%%%%%%%%%%%%%%%%%%%%%%
\subsection{Deriving the Bias--Variance Decomposition}
\label{subsec:bias-variance-decomposition}
%%%%%%%%%%%%%%%%%%%%%%%%%%%%%%%%%%%%%%%%%%%%%%%%%%%%%%%%%%%%

Write

\[
\widehat{\theta}-\theta
=
\left[
\widehat{\theta}
-
\mathbb{E}(\widehat{\theta})
\right]
+
\left[
\mathbb{E}(\widehat{\theta})
-
\theta
\right].
\]

Squaring and taking expectations gives

\[
\boxed{
\mathbb{E}
[
(
\widehat{\theta}-\theta
)^2
]
=
\operatorname{Var}(\widehat{\theta})
+
\operatorname{Bias}(\widehat{\theta})^2.
}
\]

The cross term vanishes because

\[
\mathbb{E}
[
\widehat{\theta}
-
\mathbb{E}(\widehat{\theta})
]
=
0.
\]

%%%%%%%%%%%%%%%%%%%%%%%%%%%%%%%%%%%%%%%%%%%%%%%%%%%%%%%%%%%%
\subsection{Consistency}
\label{subsec:estimator-consistency}
%%%%%%%%%%%%%%%%%%%%%%%%%%%%%%%%%%%%%%%%%%%%%%%%%%%%%%%%%%%%

\begin{definitionbox}{Consistency}
An estimator sequence \(\widehat{\theta}_n\) is consistent for
\(\theta\) if

\[
\boxed{
\widehat{\theta}_n
\xrightarrow{P}
\theta.
}
\]
\end{definitionbox}

Thus for every \(\varepsilon>0\),

\[
\boxed{
P
(
|
\widehat{\theta}_n-\theta
|>\varepsilon
)
\to0.
}
\]

Consistency is a large-sample property.

%%%%%%%%%%%%%%%%%%%%%%%%%%%%%%%%%%%%%%%%%%%%%%%%%%%%%%%%%%%%
\subsection{Consistency of the Sample Mean}
\label{subsec:consistency-sample-mean}
%%%%%%%%%%%%%%%%%%%%%%%%%%%%%%%%%%%%%%%%%%%%%%%%%%%%%%%%%%%%

If

\[
X_i
\]

are i.i.d. with finite mean

\[
\mu,
\]

the law of large numbers implies

\[
\boxed{
\overline{X}_n
\xrightarrow{P}
\mu.
}
\]

Therefore the sample mean is consistent for the population mean.

%%%%%%%%%%%%%%%%%%%%%%%%%%%%%%%%%%%%%%%%%%%%%%%%%%%%%%%%%%%%
\subsection{Mean-Square Consistency}
\label{subsec:mean-square-consistency}
%%%%%%%%%%%%%%%%%%%%%%%%%%%%%%%%%%%%%%%%%%%%%%%%%%%%%%%%%%%%

A stronger condition is

\[
\boxed{
\mathbb{E}
[
(
\widehat{\theta}_n-\theta
)^2
]
\to0.
}
\]

This is convergence in mean square.

If

\[
MSE(\widehat{\theta}_n)\to0,
\]

then

\[
\widehat{\theta}_n
\xrightarrow{P}
\theta.
\]

%%%%%%%%%%%%%%%%%%%%%%%%%%%%%%%%%%%%%%%%%%%%%%%%%%%%%%%%%%%%
\subsection{Asymptotic Normality}
\label{subsec:asymptotic-normality-estimators}
%%%%%%%%%%%%%%%%%%%%%%%%%%%%%%%%%%%%%%%%%%%%%%%%%%%%%%%%%%%%

An estimator is often asymptotically normal if

\[
\boxed{
\sqrt{n}
(
\widehat{\theta}_n-\theta
)
\xrightarrow{d}
N(0,V(\theta)).
}
\]

This implies approximately

\[
\boxed{
\widehat{\theta}_n
\approx
N
\left(
\theta,
\frac{
V(\theta)
}{n}
\right)
}
\]

for large \(n\).

Asymptotic normality is the foundation of many standard errors and
confidence intervals.

%%%%%%%%%%%%%%%%%%%%%%%%%%%%%%%%%%%%%%%%%%%%%%%%%%%%%%%%%%%%
\subsection{Method of Moments}
\label{subsec:method-of-moments}
%%%%%%%%%%%%%%%%%%%%%%%%%%%%%%%%%%%%%%%%%%%%%%%%%%%%%%%%%%%%

The method of moments estimates parameters by matching sample moments to
population moments.

Suppose

\[
\theta
\]

is determined by

\[
\mathbb{E}_\theta[X].
\]

Set

\[
\boxed{
\overline{X}
=
\mathbb{E}_\theta[X]
}
\]

and solve for \(\theta\).

For several parameters, match several moments.

%%%%%%%%%%%%%%%%%%%%%%%%%%%%%%%%%%%%%%%%%%%%%%%%%%%%%%%%%%%%
\subsection{General Method of Moments}
\label{subsec:general-method-moments}
%%%%%%%%%%%%%%%%%%%%%%%%%%%%%%%%%%%%%%%%%%%%%%%%%%%%%%%%%%%%

If there are \(k\) unknown parameters,

\[
\theta_1,\ldots,\theta_k,
\]

one may solve

\[
\boxed{
\frac1n
\sum_{i=1}^{n}
X_i^r
=
\mathbb{E}_\theta[X^r],
\qquad
r=1,\ldots,k.
}
\]

The resulting solution is called the method-of-moments estimator.

%%%%%%%%%%%%%%%%%%%%%%%%%%%%%%%%%%%%%%%%%%%%%%%%%%%%%%%%%%%%
\subsection{Worked Example: Method of Moments for Exponential Data}
\label{subsec:worked-mom-exponential}
%%%%%%%%%%%%%%%%%%%%%%%%%%%%%%%%%%%%%%%%%%%%%%%%%%%%%%%%%%%%

\begin{workedexamplebox}{Estimating an Exponential Rate}

Suppose

\[
X_i
\sim
\operatorname{Exponential}(\lambda),
\]

with

\[
\mathbb{E}[X]
=
\frac1\lambda.
\]

Set the first sample moment equal to the population moment:

\[
\overline{X}
=
\frac1\lambda.
\]

Solving gives

\begin{answerbox}
\[
\boxed{
\widehat{\lambda}_{MM}
=
\frac1{\overline{X}}.
}
\]
\end{answerbox}

\end{workedexamplebox}

%%%%%%%%%%%%%%%%%%%%%%%%%%%%%%%%%%%%%%%%%%%%%%%%%%%%%%%%%%%%
\subsection{Maximum Likelihood Estimation}
\label{subsec:maximum-likelihood-estimation-mathstat}
%%%%%%%%%%%%%%%%%%%%%%%%%%%%%%%%%%%%%%%%%%%%%%%%%%%%%%%%%%%%

Suppose observations have joint density or mass function

\[
f_\theta(x_1,\ldots,x_n).
\]

The likelihood function is

\[
\boxed{
L(\theta;x)
=
f_\theta(x_1,\ldots,x_n),
}
\]

viewed as a function of \(\theta\) for fixed observed data \(x\).

%%%%%%%%%%%%%%%%%%%%%%%%%%%%%%%%%%%%%%%%%%%%%%%%%%%%%%%%%%%%
\subsection{Maximum Likelihood Estimator}
\label{subsec:maximum-likelihood-estimator}
%%%%%%%%%%%%%%%%%%%%%%%%%%%%%%%%%%%%%%%%%%%%%%%%%%%%%%%%%%%%

The maximum likelihood estimator, abbreviated MLE, is

\[
\boxed{
\widehat{\theta}_{MLE}
=
\arg\max_{\theta\in\Theta}
L(\theta;x).
}
\]

Because logarithms preserve maximizers,

\[
\boxed{
\widehat{\theta}_{MLE}
=
\arg\max_{\theta\in\Theta}
\ell(\theta),
}
\]

where

\[
\ell(\theta)
=
\log L(\theta).
\]

%%%%%%%%%%%%%%%%%%%%%%%%%%%%%%%%%%%%%%%%%%%%%%%%%%%%%%%%%%%%
\subsection{Likelihood for an Independent Sample}
\label{subsec:likelihood-independent-sample}
%%%%%%%%%%%%%%%%%%%%%%%%%%%%%%%%%%%%%%%%%%%%%%%%%%%%%%%%%%%%

For independent observations,

\[
\boxed{
L(\theta)
=
\prod_{i=1}^{n}
f_\theta(X_i).
}
\]

Therefore,

\[
\boxed{
\ell(\theta)
=
\sum_{i=1}^{n}
\log f_\theta(X_i).
}
\]

The log-likelihood turns products into sums and is usually easier to
differentiate.

%%%%%%%%%%%%%%%%%%%%%%%%%%%%%%%%%%%%%%%%%%%%%%%%%%%%%%%%%%%%
\subsection{Worked Example: Bernoulli MLE}
\label{subsec:worked-bernoulli-mle}
%%%%%%%%%%%%%%%%%%%%%%%%%%%%%%%%%%%%%%%%%%%%%%%%%%%%%%%%%%%%

\begin{workedexamplebox}{Estimating a Bernoulli Probability}

Suppose

\[
X_i
\overset{\text{i.i.d.}}{\sim}
\operatorname{Bernoulli}(p).
\]

Then

\[
L(p)
=
\prod_{i=1}^{n}
p^{X_i}
(1-p)^{1-X_i}.
\]

Let

\[
S
=
\sum_{i=1}^{n}X_i.
\]

Then

\[
L(p)
=
p^S
(1-p)^{n-S}.
\]

The log-likelihood is

\[
\ell(p)
=
S\log p
+
(n-S)\log(1-p).
\]

Differentiate:

\[
\ell'(p)
=
\frac{S}{p}
-
\frac{n-S}{1-p}.
\]

Set equal to zero:

\[
\frac{S}{p}
=
\frac{n-S}{1-p}.
\]

Solving,

\begin{answerbox}
\[
\boxed{
\widehat{p}_{MLE}
=
\frac{S}{n}
=
\overline{X}.
}
\]
\end{answerbox}

\end{workedexamplebox}

%%%%%%%%%%%%%%%%%%%%%%%%%%%%%%%%%%%%%%%%%%%%%%%%%%%%%%%%%%%%
\subsection{Worked Example: Normal Mean MLE}
\label{subsec:worked-normal-mean-mle}
%%%%%%%%%%%%%%%%%%%%%%%%%%%%%%%%%%%%%%%%%%%%%%%%%%%%%%%%%%%%

\begin{workedexamplebox}{Normal Mean with Known Variance}

Suppose

\[
X_i
\overset{\text{i.i.d.}}{\sim}
N(\mu,\sigma^2),
\]

where \(\sigma^2\) is known.

Ignoring constants independent of \(\mu\),

\[
\ell(\mu)
=
-
\frac1{2\sigma^2}
\sum_{i=1}^{n}
(X_i-\mu)^2.
\]

Maximizing the likelihood is equivalent to minimizing

\[
\sum
(X_i-\mu)^2.
\]

Differentiating gives

\[
\sum
(X_i-\mu)=0.
\]

Therefore,

\begin{answerbox}
\[
\boxed{
\widehat{\mu}_{MLE}
=
\overline{X}.
}
\]
\end{answerbox}

\end{workedexamplebox}

%%%%%%%%%%%%%%%%%%%%%%%%%%%%%%%%%%%%%%%%%%%%%%%%%%%%%%%%%%%%
\subsection{MLE of the Normal Variance}
\label{subsec:normal-variance-mle}
%%%%%%%%%%%%%%%%%%%%%%%%%%%%%%%%%%%%%%%%%%%%%%%%%%%%%%%%%%%%

For

\[
X_i
\overset{\text{i.i.d.}}{\sim}
N(\mu,\sigma^2)
\]

with both parameters unknown, the MLEs are

\[
\boxed{
\widehat{\mu}
=
\overline{X}
}
\]

and

\[
\boxed{
\widehat{\sigma}^2_{MLE}
=
\frac1n
\sum_{i=1}^{n}
(
X_i-\overline{X}
)^2.
}
\]

Notice that the MLE uses denominator

\[
n,
\]

not \(n-1\).

%%%%%%%%%%%%%%%%%%%%%%%%%%%%%%%%%%%%%%%%%%%%%%%%%%%%%%%%%%%%
\subsection{MLE Bias Need Not Be Zero}
\label{subsec:mle-bias}
%%%%%%%%%%%%%%%%%%%%%%%%%%%%%%%%%%%%%%%%%%%%%%%%%%%%%%%%%%%%

The normal variance MLE satisfies

\[
\boxed{
\mathbb{E}
[
\widehat{\sigma}^2_{MLE}
]
=
\frac{
n-1
}{n}
\sigma^2.
}
\]

Thus it is biased downward.

The unbiased sample variance is

\[
\boxed{
S^2
=
\frac1{n-1}
\sum
(
X_i-\overline{X}
)^2.
}
\]

Maximum likelihood does not automatically produce unbiased estimators.

%%%%%%%%%%%%%%%%%%%%%%%%%%%%%%%%%%%%%%%%%%%%%%%%%%%%%%%%%%%%
\subsection{Invariance Property of Maximum Likelihood}
\label{subsec:mle-invariance}
%%%%%%%%%%%%%%%%%%%%%%%%%%%%%%%%%%%%%%%%%%%%%%%%%%%%%%%%%%%%

If

\[
\widehat{\theta}
\]

is the MLE of \(\theta\), then under appropriate conditions the MLE of

\[
g(\theta)
\]

is

\[
\boxed{
g(\widehat{\theta}).
}
\]

For example, if

\[
\widehat{\lambda}
\]

is the MLE of an exponential rate, then the MLE of the mean lifetime

\[
\frac1\lambda
\]

is

\[
\boxed{
\frac1{\widehat{\lambda}}.
}
\]

%%%%%%%%%%%%%%%%%%%%%%%%%%%%%%%%%%%%%%%%%%%%%%%%%%%%%%%%%%%%
\subsection{Score Function}
\label{subsec:score-function}
%%%%%%%%%%%%%%%%%%%%%%%%%%%%%%%%%%%%%%%%%%%%%%%%%%%%%%%%%%%%

\begin{definitionbox}{Score Function}
The score function is the derivative of the log-likelihood:

\[
\boxed{
U(\theta)
=
\frac{
\partial
}{
\partial\theta
}
\ell(\theta).
}
\]
\end{definitionbox}

For vector parameters,

\[
\boxed{
U(\boldsymbol{\theta})
=
\nabla_{\boldsymbol{\theta}}
\ell(\boldsymbol{\theta}).
}
\]

An interior MLE often satisfies

\[
\boxed{
U(\widehat{\theta})=0.
}
\]

%%%%%%%%%%%%%%%%%%%%%%%%%%%%%%%%%%%%%%%%%%%%%%%%%%%%%%%%%%%%
\subsection{Expected Score}
\label{subsec:expected-score}
%%%%%%%%%%%%%%%%%%%%%%%%%%%%%%%%%%%%%%%%%%%%%%%%%%%%%%%%%%%%

Under standard regularity conditions,

\[
\boxed{
\mathbb{E}_\theta
[
U(\theta)
]
=
0.
}
\]

This follows by differentiating

\[
\int
f_\theta(x)\,dx
=
1
\]

with respect to \(\theta\).

%%%%%%%%%%%%%%%%%%%%%%%%%%%%%%%%%%%%%%%%%%%%%%%%%%%%%%%%%%%%
\subsection{Fisher Information}
\label{subsec:fisher-information}
%%%%%%%%%%%%%%%%%%%%%%%%%%%%%%%%%%%%%%%%%%%%%%%%%%%%%%%%%%%%

\begin{definitionbox}{Fisher Information}
For a scalar parameter, Fisher information is

\[
\boxed{
I(\theta)
=
\mathbb{E}_\theta
[
U(\theta)^2
].
}
\]
\end{definitionbox}

Under regularity conditions,

\[
\boxed{
I(\theta)
=
-
\mathbb{E}_\theta
\left[
\frac{
\partial^2
}{
\partial\theta^2
}
\ell(\theta)
\right].
}
\]

Greater Fisher information corresponds to greater local sensitivity of
the distribution to changes in the parameter.

%%%%%%%%%%%%%%%%%%%%%%%%%%%%%%%%%%%%%%%%%%%%%%%%%%%%%%%%%%%%
\subsection{Information in an Independent Sample}
\label{subsec:fisher-information-independent-sample}
%%%%%%%%%%%%%%%%%%%%%%%%%%%%%%%%%%%%%%%%%%%%%%%%%%%%%%%%%%%%

For \(n\) independent identically distributed observations,

\[
\boxed{
I_n(\theta)
=
nI_1(\theta).
}
\]

Thus information typically grows linearly with sample size.

This is one mathematical reason estimator standard errors often decrease
at the rate

\[
\boxed{
n^{-1/2}.
}
\]

%%%%%%%%%%%%%%%%%%%%%%%%%%%%%%%%%%%%%%%%%%%%%%%%%%%%%%%%%%%%
\subsection{Worked Example: Bernoulli Fisher Information}
\label{subsec:worked-bernoulli-information}
%%%%%%%%%%%%%%%%%%%%%%%%%%%%%%%%%%%%%%%%%%%%%%%%%%%%%%%%%%%%

\begin{workedexamplebox}{Bernoulli Information}

For one Bernoulli observation,

\[
\ell(p)
=
X\log p
+
(1-X)\log(1-p).
\]

The score is

\[
U(p)
=
\frac{X}{p}
-
\frac{1-X}{1-p}.
\]

The second derivative is

\[
\ell''(p)
=
-
\frac{X}{p^2}
-
\frac{1-X}{(1-p)^2}.
\]

Taking negative expectation,

\[
I_1(p)
=
\frac1p
+
\frac1{1-p}.
\]

Thus,

\begin{answerbox}
\[
\boxed{
I_1(p)
=
\frac1{
p(1-p)
}.
}
\]

For \(n\) observations,

\[
\boxed{
I_n(p)
=
\frac{
n
}{
p(1-p)
}.
}
\]
\end{answerbox}

\end{workedexamplebox}

%%%%%%%%%%%%%%%%%%%%%%%%%%%%%%%%%%%%%%%%%%%%%%%%%%%%%%%%%%%%
\subsection{Cramér--Rao Lower Bound}
\label{subsec:cramer-rao-lower-bound}
%%%%%%%%%%%%%%%%%%%%%%%%%%%%%%%%%%%%%%%%%%%%%%%%%%%%%%%%%%%%

\begin{theorem}[Cramér--Rao Lower Bound]
Under suitable regularity conditions, if \(\widehat{\theta}\) is an
unbiased estimator of scalar parameter \(\theta\), then

\[
\boxed{
\operatorname{Var}_\theta
(
\widehat{\theta}
)
\geq
\frac1{
I_n(\theta)
}.
}
\]
\end{theorem}

Thus Fisher information provides a lower bound on the variance of
unbiased estimators.

%%%%%%%%%%%%%%%%%%%%%%%%%%%%%%%%%%%%%%%%%%%%%%%%%%%%%%%%%%%%
\subsection{Efficiency}
\label{subsec:estimator-efficiency}
%%%%%%%%%%%%%%%%%%%%%%%%%%%%%%%%%%%%%%%%%%%%%%%%%%%%%%%%%%%%

An unbiased estimator attaining the Cramér--Rao lower bound is called
efficient under the corresponding regular model.

For example, if

\[
X_i
\sim
N(\mu,\sigma^2)
\]

with known \(\sigma^2\),

\[
I_n(\mu)
=
\frac{n}{\sigma^2}.
\]

The bound is

\[
\frac{\sigma^2}{n}.
\]

Since

\[
\operatorname{Var}(\overline{X})
=
\frac{\sigma^2}{n},
\]

the sample mean attains the bound.

%%%%%%%%%%%%%%%%%%%%%%%%%%%%%%%%%%%%%%%%%%%%%%%%%%%%%%%%%%%%
\subsection{Sufficient Statistics}
\label{subsec:sufficient-statistics}
%%%%%%%%%%%%%%%%%%%%%%%%%%%%%%%%%%%%%%%%%%%%%%%%%%%%%%%%%%%%

\begin{definitionbox}{Sufficient Statistic}
A statistic \(T(X)\) is sufficient for parameter \(\theta\) if, once
\(T\) is known, the remaining sample information contains no additional
information about \(\theta\).
\end{definitionbox}

Formally, the conditional distribution of the full data given \(T\) does
not depend on \(\theta\).

%%%%%%%%%%%%%%%%%%%%%%%%%%%%%%%%%%%%%%%%%%%%%%%%%%%%%%%%%%%%
\subsection{Factorization Theorem}
\label{subsec:factorization-theorem}
%%%%%%%%%%%%%%%%%%%%%%%%%%%%%%%%%%%%%%%%%%%%%%%%%%%%%%%%%%%%

\begin{theorem}[Fisher--Neyman Factorization Theorem]
A statistic \(T(X)\) is sufficient for \(\theta\) if the joint density or
mass function can be written

\[
\boxed{
f_\theta(x)
=
g_\theta
(
T(x)
)
h(x),
}
\]

where \(h(x)\) does not depend on \(\theta\).
\end{theorem}

This provides a practical method for identifying sufficient statistics.

%%%%%%%%%%%%%%%%%%%%%%%%%%%%%%%%%%%%%%%%%%%%%%%%%%%%%%%%%%%%
\subsection{Worked Example: Bernoulli Sufficiency}
\label{subsec:worked-bernoulli-sufficiency}
%%%%%%%%%%%%%%%%%%%%%%%%%%%%%%%%%%%%%%%%%%%%%%%%%%%%%%%%%%%%

\begin{workedexamplebox}{Sufficient Statistic for Bernoulli Data}

Suppose

\[
X_i
\overset{\text{i.i.d.}}{\sim}
\operatorname{Bernoulli}(p).
\]

The joint mass function is

\[
f_p(x)
=
p^{\sum x_i}
(1-p)^{n-\sum x_i}.
\]

This depends on the sample only through

\[
T
=
\sum_{i=1}^{n}X_i.
\]

Therefore, by the factorization theorem,

\begin{answerbox}
\[
\boxed{
T
=
\sum_{i=1}^{n}X_i
}
\]

is sufficient for \(p\).
\end{answerbox}

\end{workedexamplebox}

%%%%%%%%%%%%%%%%%%%%%%%%%%%%%%%%%%%%%%%%%%%%%%%%%%%%%%%%%%%%
\subsection{Worked Example: Normal Sufficiency}
\label{subsec:worked-normal-sufficiency}
%%%%%%%%%%%%%%%%%%%%%%%%%%%%%%%%%%%%%%%%%%%%%%%%%%%%%%%%%%%%

Suppose

\[
X_i
\overset{\text{i.i.d.}}{\sim}
N(\mu,\sigma^2),
\]

with both parameters unknown.

The likelihood depends on the sample through

\[
\sum X_i
\]

and

\[
\sum X_i^2.
\]

Thus

\[
\boxed{
\left(
\sum X_i,
\sum X_i^2
\right)
}
\]

is sufficient for

\[
(\mu,\sigma^2).
\]

Equivalent one-to-one forms include

\[
\boxed{
(
\overline{X},
S^2
).
}
\]

%%%%%%%%%%%%%%%%%%%%%%%%%%%%%%%%%%%%%%%%%%%%%%%%%%%%%%%%%%%%
\subsection{Minimal Sufficiency}
\label{subsec:minimal-sufficiency}
%%%%%%%%%%%%%%%%%%%%%%%%%%%%%%%%%%%%%%%%%%%%%%%%%%%%%%%%%%%%

A sufficient statistic may contain redundant information.

A minimal sufficient statistic compresses the sample as far as possible
while retaining all parameter information in a precise statistical
sense.

A statistic \(T\) is minimal sufficient if every other sufficient
statistic determines \(T\), up to the relevant equivalence.

%%%%%%%%%%%%%%%%%%%%%%%%%%%%%%%%%%%%%%%%%%%%%%%%%%%%%%%%%%%%
\subsection{Ancillary Statistics}
\label{subsec:ancillary-statistics}
%%%%%%%%%%%%%%%%%%%%%%%%%%%%%%%%%%%%%%%%%%%%%%%%%%%%%%%%%%%%

\begin{definitionbox}{Ancillary Statistic}
A statistic is ancillary if its distribution does not depend on the
unknown parameter.
\end{definitionbox}

Ancillary statistics contain no direct information about the parameter
but can describe sample geometry or precision.

Conditioning on ancillary information plays an important role in advanced
frequentist inference.

%%%%%%%%%%%%%%%%%%%%%%%%%%%%%%%%%%%%%%%%%%%%%%%%%%%%%%%%%%%%
\subsection{Complete Statistics}
\label{subsec:complete-statistics}
%%%%%%%%%%%%%%%%%%%%%%%%%%%%%%%%%%%%%%%%%%%%%%%%%%%%%%%%%%%%

\begin{definitionbox}{Complete Statistic}
A statistic \(T\) is complete for a family if

\[
\boxed{
\mathbb{E}_\theta[g(T)]=0
\text{ for all }\theta
}
\]

implies

\[
\boxed{
P_\theta(g(T)=0)=1
\text{ for all }\theta.
}
\]
\end{definitionbox}

Completeness is a strong uniqueness property that becomes especially
powerful when combined with sufficiency.

%%%%%%%%%%%%%%%%%%%%%%%%%%%%%%%%%%%%%%%%%%%%%%%%%%%%%%%%%%%%
\subsection{Rao--Blackwell Theorem}
\label{subsec:rao-blackwell-theorem}
%%%%%%%%%%%%%%%%%%%%%%%%%%%%%%%%%%%%%%%%%%%%%%%%%%%%%%%%%%%%

\begin{theorem}[Rao--Blackwell]
Let \(U\) be an estimator of \(\theta\), and let \(T\) be sufficient for
\(\theta\). Define

\[
\boxed{
U^*
=
\mathbb{E}[U\mid T].
}
\]

Then \(U^*\) has the same expectation as \(U\) and satisfies

\[
\boxed{
\operatorname{Var}(U^*)
\leq
\operatorname{Var}(U).
}
\]
\end{theorem}

Conditioning on a sufficient statistic can therefore improve an estimator.

%%%%%%%%%%%%%%%%%%%%%%%%%%%%%%%%%%%%%%%%%%%%%%%%%%%%%%%%%%%%
\subsection{Why Rao--Blackwell Improves Variance}
\label{subsec:rao-blackwell-variance}
%%%%%%%%%%%%%%%%%%%%%%%%%%%%%%%%%%%%%%%%%%%%%%%%%%%%%%%%%%%%

By the law of total variance,

\[
\operatorname{Var}(U)
=
\operatorname{Var}
(
\mathbb{E}[U\mid T]
)
+
\mathbb{E}
[
\operatorname{Var}(U\mid T)
].
\]

Therefore,

\[
\boxed{
\operatorname{Var}
(
\mathbb{E}[U\mid T]
)
\leq
\operatorname{Var}(U).
}
\]

The difference is the expected conditional variance.

%%%%%%%%%%%%%%%%%%%%%%%%%%%%%%%%%%%%%%%%%%%%%%%%%%%%%%%%%%%%
\subsection{Lehmann--Scheffé Theorem}
\label{subsec:lehmann-scheffe-theorem}
%%%%%%%%%%%%%%%%%%%%%%%%%%%%%%%%%%%%%%%%%%%%%%%%%%%%%%%%%%%%

\begin{theorem}[Lehmann--Scheffé]
If \(T\) is complete and sufficient for \(\theta\), then any unbiased
estimator that is a function of \(T\) is the unique minimum-variance
unbiased estimator of its expectation.
\end{theorem}

Thus complete sufficient statistics identify optimal unbiased estimators.

%%%%%%%%%%%%%%%%%%%%%%%%%%%%%%%%%%%%%%%%%%%%%%%%%%%%%%%%%%%%
\subsection{UMVU Estimators}
\label{subsec:umvu-estimators}
%%%%%%%%%%%%%%%%%%%%%%%%%%%%%%%%%%%%%%%%%%%%%%%%%%%%%%%%%%%%

UMVU stands for uniformly minimum-variance unbiased.

An estimator \(\widehat{\tau}\) of \(\tau(\theta)\) is UMVU if it is
unbiased and has variance no larger than any other unbiased estimator for
every \(\theta\).

The Lehmann--Scheffé theorem provides a major method for proving UMVU
optimality.

%%%%%%%%%%%%%%%%%%%%%%%%%%%%%%%%%%%%%%%%%%%%%%%%%%%%%%%%%%%%
\subsection{Exponential Families}
\label{subsec:exponential-families}
%%%%%%%%%%%%%%%%%%%%%%%%%%%%%%%%%%%%%%%%%%%%%%%%%%%%%%%%%%%%

A distribution belongs to a \(k\)-parameter exponential family if its
density or mass function can be written

\[
\boxed{
f_\theta(x)
=
h(x)
\exp
\left[
\sum_{j=1}^{k}
\eta_j(\theta)
T_j(x)
-
A(\theta)
\right].
}
\]

The functions

\[
T_1,\ldots,T_k
\]

are natural sufficient statistics.

%%%%%%%%%%%%%%%%%%%%%%%%%%%%%%%%%%%%%%%%%%%%%%%%%%%%%%%%%%%%
\subsection{Natural Exponential-Family Form}
\label{subsec:natural-exponential-family}
%%%%%%%%%%%%%%%%%%%%%%%%%%%%%%%%%%%%%%%%%%%%%%%%%%%%%%%%%%%%

A canonical form is

\[
\boxed{
f_\eta(x)
=
h(x)
\exp
\left[
\eta^T T(x)
-
A(\eta)
\right].
}
\]

Here,

\[
\eta
\]

is the natural parameter,

\[
T(x)
\]

is the natural statistic, and

\[
A(\eta)
\]

is the log-partition function.

%%%%%%%%%%%%%%%%%%%%%%%%%%%%%%%%%%%%%%%%%%%%%%%%%%%%%%%%%%%%
\subsection{Moments from the Log-Partition Function}
\label{subsec:log-partition-moments}
%%%%%%%%%%%%%%%%%%%%%%%%%%%%%%%%%%%%%%%%%%%%%%%%%%%%%%%%%%%%

Under regular conditions,

\[
\boxed{
\nabla A(\eta)
=
\mathbb{E}_\eta[T(X)].
}
\]

Also,

\[
\boxed{
\nabla^2 A(\eta)
=
\operatorname{Cov}_\eta
(
T(X)
).
}
\]

Thus the geometry of the log-partition function encodes moments and
information.

%%%%%%%%%%%%%%%%%%%%%%%%%%%%%%%%%%%%%%%%%%%%%%%%%%%%%%%%%%%%
\subsection{Examples of Exponential Families}
\label{subsec:exponential-family-examples}
%%%%%%%%%%%%%%%%%%%%%%%%%%%%%%%%%%%%%%%%%%%%%%%%%%%%%%%%%%%%

Important exponential-family distributions include:

\[
\boxed{
\text{Bernoulli},
}
\]

\[
\boxed{
\text{binomial},
}
\]

\[
\boxed{
\text{Poisson},
}
\]

\[
\boxed{
\text{normal},
}
\]

\[
\boxed{
\text{gamma under suitable parameterizations},
}
\]

and many others.

This structure explains why similar inference methods recur across many
models.

%%%%%%%%%%%%%%%%%%%%%%%%%%%%%%%%%%%%%%%%%%%%%%%%%%%%%%%%%%%%
\subsection{Pivotal Quantities}
\label{subsec:pivotal-quantities}
%%%%%%%%%%%%%%%%%%%%%%%%%%%%%%%%%%%%%%%%%%%%%%%%%%%%%%%%%%%%

\begin{definitionbox}{Pivotal Quantity}
A pivotal quantity is a function of the data and unknown parameter whose
sampling distribution does not depend on unknown parameters.
\end{definitionbox}

For example, if

\[
X_i
\sim
N(\mu,\sigma^2)
\]

with known \(\sigma\),

\[
\boxed{
Z
=
\frac{
\overline{X}-\mu
}{
\sigma/\sqrt{n}
}
\sim
N(0,1).
}
\]

The distribution does not depend on \(\mu\).

%%%%%%%%%%%%%%%%%%%%%%%%%%%%%%%%%%%%%%%%%%%%%%%%%%%%%%%%%%%%
\subsection{Pivots and Confidence Intervals}
\label{subsec:pivots-confidence-intervals}
%%%%%%%%%%%%%%%%%%%%%%%%%%%%%%%%%%%%%%%%%%%%%%%%%%%%%%%%%%%%

Suppose

\[
Q(X,\theta)
\]

is a pivot and

\[
P
(
a
\leq
Q(X,\theta)
\leq
b
)
=
1-\alpha.
\]

Solve the inequalities for \(\theta\).

This produces a confidence interval with exact coverage under the model.

%%%%%%%%%%%%%%%%%%%%%%%%%%%%%%%%%%%%%%%%%%%%%%%%%%%%%%%%%%%%
\subsection{Normal Mean Pivot}
\label{subsec:normal-mean-pivot}
%%%%%%%%%%%%%%%%%%%%%%%%%%%%%%%%%%%%%%%%%%%%%%%%%%%%%%%%%%%%

If

\[
X_i
\overset{\text{i.i.d.}}{\sim}
N(\mu,\sigma^2)
\]

with unknown \(\sigma\), then

\[
\boxed{
T
=
\frac{
\overline{X}-\mu
}{
S/\sqrt{n}
}
\sim
t_{n-1}.
}
\]

This pivot produces the exact classical confidence interval

\[
\boxed{
\overline{X}
\pm
t^*_{n-1}
\frac{S}{\sqrt{n}}.
}
\]

%%%%%%%%%%%%%%%%%%%%%%%%%%%%%%%%%%%%%%%%%%%%%%%%%%%%%%%%%%%%
\subsection{Normal Variance Pivot}
\label{subsec:normal-variance-pivot}
%%%%%%%%%%%%%%%%%%%%%%%%%%%%%%%%%%%%%%%%%%%%%%%%%%%%%%%%%%%%

For normal data,

\[
\boxed{
\frac{
(n-1)S^2
}{
\sigma^2
}
\sim
\chi^2_{n-1}.
}
\]

Therefore exact intervals for \(\sigma^2\) can be obtained by inverting
chi-square probability bounds.

%%%%%%%%%%%%%%%%%%%%%%%%%%%%%%%%%%%%%%%%%%%%%%%%%%%%%%%%%%%%
\subsection{Confidence Sets by Test Inversion}
\label{subsec:confidence-set-test-inversion}
%%%%%%%%%%%%%%%%%%%%%%%%%%%%%%%%%%%%%%%%%%%%%%%%%%%%%%%%%%%%

Suppose that for each candidate value \(\theta_0\), a level-\(\alpha\)
test is available.

The corresponding \(100(1-\alpha)\%\) confidence set contains all values
\(\theta_0\) not rejected by the test.

Thus

\[
\boxed{
\text{confidence set}
=
\{
\theta_0:
H_0:\theta=\theta_0
\text{ is not rejected}
\}.
}
\]

This establishes a general test--confidence-set duality.

%%%%%%%%%%%%%%%%%%%%%%%%%%%%%%%%%%%%%%%%%%%%%%%%%%%%%%%%%%%%
\subsection{Neyman--Pearson Lemma}
\label{subsec:neyman-pearson-mathstat}
%%%%%%%%%%%%%%%%%%%%%%%%%%%%%%%%%%%%%%%%%%%%%%%%%%%%%%%%%%%%

Suppose

\[
H_0:\theta=\theta_0
\]

and

\[
H_1:\theta=\theta_1
\]

are both simple.

\begin{theorem}[Neyman--Pearson Lemma]
Among all tests of size at most \(\alpha\), a most powerful test rejects
for sufficiently large values of

\[
\boxed{
\frac{
f_{\theta_1}(x)
}{
f_{\theta_0}(x)
}.
}
\]
\end{theorem}

Thus likelihood ratios arise naturally from optimal testing.

%%%%%%%%%%%%%%%%%%%%%%%%%%%%%%%%%%%%%%%%%%%%%%%%%%%%%%%%%%%%
\subsection{Power Function}
\label{subsec:power-function-mathstat}
%%%%%%%%%%%%%%%%%%%%%%%%%%%%%%%%%%%%%%%%%%%%%%%%%%%%%%%%%%%%

For a test with rejection region \(R\), the power function is

\[
\boxed{
\pi(\theta)
=
P_\theta(X\in R).
}
\]

For null parameter values, \(\pi(\theta)\) describes Type I rejection
probabilities.

For alternatives, it describes the probability of detecting departures
from the null.

%%%%%%%%%%%%%%%%%%%%%%%%%%%%%%%%%%%%%%%%%%%%%%%%%%%%%%%%%%%%
\subsection{Most Powerful Tests}
\label{subsec:most-powerful-mathstat}
%%%%%%%%%%%%%%%%%%%%%%%%%%%%%%%%%%%%%%%%%%%%%%%%%%%%%%%%%%%%

A test is most powerful against a specified alternative \(\theta_1\) if
its power at \(\theta_1\) is at least as large as that of every other
test with the same size constraint.

This is a pointwise optimality statement.

It does not automatically guarantee optimality against all alternatives.

%%%%%%%%%%%%%%%%%%%%%%%%%%%%%%%%%%%%%%%%%%%%%%%%%%%%%%%%%%%%
\subsection{Uniformly Most Powerful Tests}
\label{subsec:ump-mathstat}
%%%%%%%%%%%%%%%%%%%%%%%%%%%%%%%%%%%%%%%%%%%%%%%%%%%%%%%%%%%%

A uniformly most powerful test, abbreviated UMP, maximizes power
simultaneously for every parameter value in a specified alternative
space.

Such tests do not always exist.

They are especially important in one-parameter exponential families with
monotone likelihood-ratio structure.

%%%%%%%%%%%%%%%%%%%%%%%%%%%%%%%%%%%%%%%%%%%%%%%%%%%%%%%%%%%%
\subsection{Monotone Likelihood Ratio}
\label{subsec:monotone-likelihood-ratio}
%%%%%%%%%%%%%%%%%%%%%%%%%%%%%%%%%%%%%%%%%%%%%%%%%%%%%%%%%%%%

A family has a monotone likelihood ratio in statistic \(T(X)\) if for
\(\theta_2>\theta_1\),

\[
\boxed{
\frac{
f_{\theta_2}(x)
}{
f_{\theta_1}(x)
}
}
\]

is a nondecreasing function of \(T(x)\).

This structure often implies that large values of \(T\) provide evidence
for larger parameter values.

%%%%%%%%%%%%%%%%%%%%%%%%%%%%%%%%%%%%%%%%%%%%%%%%%%%%%%%%%%%%
\subsection{Karlin--Rubin Theorem Preview}
\label{subsec:karlin-rubin-preview}
%%%%%%%%%%%%%%%%%%%%%%%%%%%%%%%%%%%%%%%%%%%%%%%%%%%%%%%%%%%%

Under appropriate monotone likelihood-ratio conditions, one-sided tests
that reject for large values of a sufficient statistic can be UMP.

This result extends Neyman--Pearson theory from simple alternatives to
certain composite alternatives.

%%%%%%%%%%%%%%%%%%%%%%%%%%%%%%%%%%%%%%%%%%%%%%%%%%%%%%%%%%%%
\subsection{Likelihood-Ratio Tests}
\label{subsec:likelihood-ratio-tests-mathstat}
%%%%%%%%%%%%%%%%%%%%%%%%%%%%%%%%%%%%%%%%%%%%%%%%%%%%%%%%%%%%

Let the full parameter space be

\[
\Theta
\]

and the null parameter space be

\[
\Theta_0.
\]

Define

\[
\boxed{
\Lambda
=
\frac{
\sup_{\theta\in\Theta_0}
L(\theta)
}{
\sup_{\theta\in\Theta}
L(\theta)
}.
}
\]

Because the numerator maximizes over a smaller set,

\[
\boxed{
0\leq\Lambda\leq1.
}
\]

Small values provide evidence against the null.

%%%%%%%%%%%%%%%%%%%%%%%%%%%%%%%%%%%%%%%%%%%%%%%%%%%%%%%%%%%%
\subsection{Likelihood-Ratio Statistic}
\label{subsec:likelihood-ratio-statistic-mathstat}
%%%%%%%%%%%%%%%%%%%%%%%%%%%%%%%%%%%%%%%%%%%%%%%%%%%%%%%%%%%%

The transformed likelihood-ratio statistic is

\[
\boxed{
-2\log\Lambda.
}
\]

This quantity is nonnegative and often has a simple asymptotic
distribution.

%%%%%%%%%%%%%%%%%%%%%%%%%%%%%%%%%%%%%%%%%%%%%%%%%%%%%%%%%%%%
\subsection{Wilks' Theorem}
\label{subsec:wilks-theorem-mathstat}
%%%%%%%%%%%%%%%%%%%%%%%%%%%%%%%%%%%%%%%%%%%%%%%%%%%%%%%%%%%%

Under standard regularity conditions, if the null imposes \(q\)
independent restrictions,

\[
\boxed{
-2\log\Lambda
\xrightarrow{d}
\chi^2_q.
}
\]

This result provides a general asymptotic likelihood-ratio test.

Regularity can fail at parameter-space boundaries or in nonidentifiable
models.

%%%%%%%%%%%%%%%%%%%%%%%%%%%%%%%%%%%%%%%%%%%%%%%%%%%%%%%%%%%%
\subsection{Wald Tests}
\label{subsec:wald-tests-mathstat}
%%%%%%%%%%%%%%%%%%%%%%%%%%%%%%%%%%%%%%%%%%%%%%%%%%%%%%%%%%%%

Suppose

\[
\widehat{\theta}
\]

is asymptotically normal:

\[
\sqrt{n}
(
\widehat{\theta}-\theta
)
\xrightarrow{d}
N(0,V(\theta)).
\]

A Wald statistic for

\[
H_0:\theta=\theta_0
\]

is

\[
\boxed{
W
=
\frac{
(
\widehat{\theta}-\theta_0
)^2
}{
\widehat{\operatorname{Var}}
(
\widehat{\theta}
)
}.
}
\]

Under the null,

\[
\boxed{
W
\approx
\chi^2_1.
}
\]

%%%%%%%%%%%%%%%%%%%%%%%%%%%%%%%%%%%%%%%%%%%%%%%%%%%%%%%%%%%%
\subsection{Score Tests}
\label{subsec:score-tests-mathstat}
%%%%%%%%%%%%%%%%%%%%%%%%%%%%%%%%%%%%%%%%%%%%%%%%%%%%%%%%%%%%

The score test evaluates the score at the null value:

\[
U(\theta_0).
\]

A scalar score statistic is

\[
\boxed{
S
=
\frac{
U(\theta_0)^2
}{
I(\theta_0)
}.
}
\]

Under standard regularity conditions,

\[
\boxed{
S
\approx
\chi^2_1.
}
\]

The score test requires fitting only the null model in many settings.

%%%%%%%%%%%%%%%%%%%%%%%%%%%%%%%%%%%%%%%%%%%%%%%%%%%%%%%%%%%%
\subsection{Wald, Score, and Likelihood Ratio}
\label{subsec:wald-score-lr-mathstat}
%%%%%%%%%%%%%%%%%%%%%%%%%%%%%%%%%%%%%%%%%%%%%%%%%%%%%%%%%%%%

In regular large samples,

\[
\boxed{
\text{Wald}
\approx
\text{Score}
\approx
\text{Likelihood Ratio}.
}
\]

Their differences vanish asymptotically under suitable local
alternatives.

In finite samples they may give noticeably different results.

%%%%%%%%%%%%%%%%%%%%%%%%%%%%%%%%%%%%%%%%%%%%%%%%%%%%%%%%%%%%
\subsection{Asymptotic Distribution of the MLE}
\label{subsec:asymptotic-mle}
%%%%%%%%%%%%%%%%%%%%%%%%%%%%%%%%%%%%%%%%%%%%%%%%%%%%%%%%%%%%

Under standard regularity conditions,

\[
\boxed{
\sqrt{n}
(
\widehat{\theta}_{MLE}-\theta
)
\xrightarrow{d}
N
\left(
0,
I_1(\theta)^{-1}
\right).
}
\]

Equivalently,

\[
\boxed{
\widehat{\theta}_{MLE}
\approx
N
\left(
\theta,
\frac1{
nI_1(\theta)
}
\right).
}
\]

Thus the MLE is asymptotically efficient in many regular models.

%%%%%%%%%%%%%%%%%%%%%%%%%%%%%%%%%%%%%%%%%%%%%%%%%%%%%%%%%%%%
\subsection{Observed Information}
\label{subsec:observed-information}
%%%%%%%%%%%%%%%%%%%%%%%%%%%%%%%%%%%%%%%%%%%%%%%%%%%%%%%%%%%%

The observed information is

\[
\boxed{
J(\theta)
=
-
\frac{
\partial^2
}{
\partial\theta^2
}
\ell(\theta).
}
\]

Evaluated at the MLE,

\[
\boxed{
J(\widehat{\theta})
}
\]

can be used to approximate the estimator covariance.

A common large-sample standard error is

\[
\boxed{
SE(\widehat{\theta})
\approx
\frac1{
\sqrt{
J(\widehat{\theta})
}
}.
}
\]

%%%%%%%%%%%%%%%%%%%%%%%%%%%%%%%%%%%%%%%%%%%%%%%%%%%%%%%%%%%%
\subsection{Delta Method}
\label{subsec:delta-method-mathstat}
%%%%%%%%%%%%%%%%%%%%%%%%%%%%%%%%%%%%%%%%%%%%%%%%%%%%%%%%%%%%

Suppose

\[
\sqrt{n}
(
\widehat{\theta}-\theta
)
\xrightarrow{d}
N(0,V)
\]

and \(g\) is differentiable at \(\theta\).

Then

\[
\boxed{
\sqrt{n}
[
g(\widehat{\theta})-g(\theta)
]
\xrightarrow{d}
N
\left(
0,
[g'(\theta)]^2V
\right).
}
\]

This is the delta method.

%%%%%%%%%%%%%%%%%%%%%%%%%%%%%%%%%%%%%%%%%%%%%%%%%%%%%%%%%%%%
\subsection{Worked Example: Delta Method}
\label{subsec:worked-delta-method}
%%%%%%%%%%%%%%%%%%%%%%%%%%%%%%%%%%%%%%%%%%%%%%%%%%%%%%%%%%%%

\begin{workedexamplebox}{Log Transformation}

Suppose

\[
\widehat{\theta}
\approx
N
\left(
\theta,
\frac{V}{n}
\right),
\qquad
\theta>0.
\]

Let

\[
g(\theta)=\log\theta.
\]

Then

\[
g'(\theta)=\frac1\theta.
\]

Therefore,

\begin{answerbox}
\[
\boxed{
\log\widehat{\theta}
\approx
N
\left(
\log\theta,
\frac{
V
}{
n\theta^2
}
\right).
}
\]
\end{answerbox}

\end{workedexamplebox}

%%%%%%%%%%%%%%%%%%%%%%%%%%%%%%%%%%%%%%%%%%%%%%%%%%%%%%%%%%%%
\subsection{Multivariate Delta Method}
\label{subsec:multivariate-delta-method-mathstat}
%%%%%%%%%%%%%%%%%%%%%%%%%%%%%%%%%%%%%%%%%%%%%%%%%%%%%%%%%%%%

If

\[
\sqrt{n}
(
\widehat{\boldsymbol{\theta}}
-
\boldsymbol{\theta}
)
\xrightarrow{d}
N
(
0,\Sigma
)
\]

and

\[
g:\mathbb{R}^k\to\mathbb{R}^m
\]

is differentiable, then

\[
\boxed{
\sqrt{n}
[
g(\widehat{\boldsymbol{\theta}})
-
g(\boldsymbol{\theta})
]
\xrightarrow{d}
N
\left(
0,
Dg(\boldsymbol{\theta})
\Sigma
Dg(\boldsymbol{\theta})^T
\right).
}
\]

Here \(Dg\) is the Jacobian matrix.

%%%%%%%%%%%%%%%%%%%%%%%%%%%%%%%%%%%%%%%%%%%%%%%%%%%%%%%%%%%%
\subsection{Slutsky's Theorem in Statistics}
\label{subsec:slutsky-statistics}
%%%%%%%%%%%%%%%%%%%%%%%%%%%%%%%%%%%%%%%%%%%%%%%%%%%%%%%%%%%%

Suppose

\[
X_n
\xrightarrow{d}
X
\]

and

\[
Y_n
\xrightarrow{P}
c.
\]

Then

\[
\boxed{
X_n+Y_n
\xrightarrow{d}
X+c,
}
\]

and

\[
\boxed{
X_nY_n
\xrightarrow{d}
cX.
}
\]

If \(c\neq0\),

\[
\boxed{
\frac{X_n}{Y_n}
\xrightarrow{d}
\frac{X}{c}.
}
\]

This justifies replacing unknown nuisance parameters by consistent
estimators in many asymptotic statistics.

%%%%%%%%%%%%%%%%%%%%%%%%%%%%%%%%%%%%%%%%%%%%%%%%%%%%%%%%%%%%
\subsection{Continuous Mapping Theorem}
\label{subsec:continuous-mapping-statistics}
%%%%%%%%%%%%%%%%%%%%%%%%%%%%%%%%%%%%%%%%%%%%%%%%%%%%%%%%%%%%

If

\[
X_n
\xrightarrow{P}
X
\]

and \(g\) is continuous, then

\[
\boxed{
g(X_n)
\xrightarrow{P}
g(X).
}
\]

Likewise, convergence in distribution is preserved under continuous
mappings.

This is fundamental for proving consistency of transformed estimators.

%%%%%%%%%%%%%%%%%%%%%%%%%%%%%%%%%%%%%%%%%%%%%%%%%%%%%%%%%%%%
\subsection{Asymptotic Relative Efficiency}
\label{subsec:asymptotic-relative-efficiency-mathstat}
%%%%%%%%%%%%%%%%%%%%%%%%%%%%%%%%%%%%%%%%%%%%%%%%%%%%%%%%%%%%

Suppose two estimators satisfy

\[
\sqrt{n}
(
\widehat{\theta}_1-\theta
)
\xrightarrow{d}
N(0,V_1)
\]

and

\[
\sqrt{n}
(
\widehat{\theta}_2-\theta
)
\xrightarrow{d}
N(0,V_2).
\]

A variance-based asymptotic efficiency comparison uses ratios such as

\[
\boxed{
\frac{
V_2
}{
V_1
}.
}
\]

Smaller asymptotic variance means greater large-sample precision.

%%%%%%%%%%%%%%%%%%%%%%%%%%%%%%%%%%%%%%%%%%%%%%%%%%%%%%%%%%%%
\subsection{Order Statistics}
\label{subsec:order-statistics-mathstat}
%%%%%%%%%%%%%%%%%%%%%%%%%%%%%%%%%%%%%%%%%%%%%%%%%%%%%%%%%%%%

For an i.i.d. sample,

\[
X_1,\ldots,X_n,
\]

the ordered values are

\[
\boxed{
X_{(1)}
\leq
X_{(2)}
\leq
\cdots
\leq
X_{(n)}.
}
\]

Order statistics are fundamental for:

\[
\boxed{
\text{quantiles},
}
\]

\[
\boxed{
\text{extremes},
}
\]

\[
\boxed{
\text{nonparametric inference},
}
\]

and

\[
\boxed{
\text{reliability}.
}
\]

%%%%%%%%%%%%%%%%%%%%%%%%%%%%%%%%%%%%%%%%%%%%%%%%%%%%%%%%%%%%
\subsection{Distribution of an Order Statistic}
\label{subsec:order-statistic-density}
%%%%%%%%%%%%%%%%%%%%%%%%%%%%%%%%%%%%%%%%%%%%%%%%%%%%%%%%%%%%

For continuous CDF \(F\) and density \(f\), the density of the
\(k\)-th order statistic is

\[
\boxed{
f_{X_{(k)}}(x)
=
\frac{
n!
}{
(k-1)!(n-k)!
}
[F(x)]^{k-1}
[1-F(x)]^{n-k}
f(x).
}
\]

This formula reflects:

\[
k-1
\]

observations below \(x\),

one observation near \(x\), and

\[
n-k
\]

observations above \(x\).

%%%%%%%%%%%%%%%%%%%%%%%%%%%%%%%%%%%%%%%%%%%%%%%%%%%%%%%%%%%%
\subsection{Distribution of the Minimum}
\label{subsec:minimum-order-statistic}
%%%%%%%%%%%%%%%%%%%%%%%%%%%%%%%%%%%%%%%%%%%%%%%%%%%%%%%%%%%%

For

\[
X_{(1)}
=
\min_iX_i,
\]

\[
P
(
X_{(1)}>x
)
=
P
(
X_1>x,\ldots,X_n>x
).
\]

Under independence,

\[
\boxed{
P
(
X_{(1)}>x
)
=
[1-F(x)]^n.
}
\]

Therefore,

\[
\boxed{
F_{X_{(1)}}(x)
=
1-
[1-F(x)]^n.
}
\]

%%%%%%%%%%%%%%%%%%%%%%%%%%%%%%%%%%%%%%%%%%%%%%%%%%%%%%%%%%%%
\subsection{Distribution of the Maximum}
\label{subsec:maximum-order-statistic}
%%%%%%%%%%%%%%%%%%%%%%%%%%%%%%%%%%%%%%%%%%%%%%%%%%%%%%%%%%%%

For

\[
X_{(n)}
=
\max_iX_i,
\]

\[
P
(
X_{(n)}\leq x
)
=
P
(
X_1\leq x,\ldots,X_n\leq x
).
\]

Thus,

\[
\boxed{
F_{X_{(n)}}(x)
=
[F(x)]^n.
}
\]

%%%%%%%%%%%%%%%%%%%%%%%%%%%%%%%%%%%%%%%%%%%%%%%%%%%%%%%%%%%%
\subsection{Probability Integral Transform and Order Statistics}
\label{subsec:pit-order-statistics}
%%%%%%%%%%%%%%%%%%%%%%%%%%%%%%%%%%%%%%%%%%%%%%%%%%%%%%%%%%%%

If \(F\) is continuous, then

\[
U_i
=
F(X_i)
\]

are i.i.d.

\[
\operatorname{Uniform}(0,1).
\]

Therefore,

\[
\boxed{
F
(
X_{(k)}
)
\sim
\operatorname{Beta}
(
k,
n-k+1
).
}
\]

This provides a powerful route to distribution theory for order
statistics.

%%%%%%%%%%%%%%%%%%%%%%%%%%%%%%%%%%%%%%%%%%%%%%%%%%%%%%%%%%%%
\subsection{Statistical Decision Theory}
\label{subsec:statistical-decision-theory}
%%%%%%%%%%%%%%%%%%%%%%%%%%%%%%%%%%%%%%%%%%%%%%%%%%%%%%%%%%%%

Statistical inference can be viewed as a decision problem.

A decision rule maps observed data into an action.

Let the action be

\[
a.
\]

The quality of the action under parameter \(\theta\) is measured by a
loss function.

%%%%%%%%%%%%%%%%%%%%%%%%%%%%%%%%%%%%%%%%%%%%%%%%%%%%%%%%%%%%
\subsection{Loss Function}
\label{subsec:loss-function}
%%%%%%%%%%%%%%%%%%%%%%%%%%%%%%%%%%%%%%%%%%%%%%%%%%%%%%%%%%%%

\begin{definitionbox}{Loss Function}
A loss function

\[
\boxed{
L(\theta,a)
}
\]

measures the cost of taking action \(a\) when the true parameter is
\(\theta\).
\end{definitionbox}

For point estimation, a common loss is squared error:

\[
\boxed{
L
(
\theta,a
)
=
(a-\theta)^2.
}
\]

%%%%%%%%%%%%%%%%%%%%%%%%%%%%%%%%%%%%%%%%%%%%%%%%%%%%%%%%%%%%
\subsection{Risk Function}
\label{subsec:risk-function}
%%%%%%%%%%%%%%%%%%%%%%%%%%%%%%%%%%%%%%%%%%%%%%%%%%%%%%%%%%%%

If decision rule \(\delta(X)\) produces action

\[
a=\delta(X),
\]

its risk is

\[
\boxed{
R
(
\theta,\delta
)
=
\mathbb{E}_\theta
[
L
(
\theta,\delta(X)
)
].
}
\]

Under squared-error loss, the risk of an estimator is exactly its mean
squared error.

%%%%%%%%%%%%%%%%%%%%%%%%%%%%%%%%%%%%%%%%%%%%%%%%%%%%%%%%%%%%
\subsection{Bayes Risk}
\label{subsec:bayes-risk}
%%%%%%%%%%%%%%%%%%%%%%%%%%%%%%%%%%%%%%%%%%%%%%%%%%%%%%%%%%%%

If a prior distribution

\[
\pi(\theta)
\]

is placed on \(\theta\), the Bayes risk is

\[
\boxed{
r
(
\pi,\delta
)
=
\int
R
(
\theta,\delta
)
\pi(\theta)\,d\theta.
}
\]

A Bayes rule minimizes this prior-averaged risk.

%%%%%%%%%%%%%%%%%%%%%%%%%%%%%%%%%%%%%%%%%%%%%%%%%%%%%%%%%%%%
\subsection{Bayes Estimator Under Squared Error}
\label{subsec:bayes-estimator-squared-error}
%%%%%%%%%%%%%%%%%%%%%%%%%%%%%%%%%%%%%%%%%%%%%%%%%%%%%%%%%%%%

Under squared-error loss,

\[
L(\theta,a)
=
(a-\theta)^2,
\]

the Bayes estimator is the posterior mean:

\[
\boxed{
\delta_B(x)
=
\mathbb{E}
[
\theta\mid X=x
].
}
\]

Under absolute-error loss, the posterior median is a Bayes rule.

Under zero-one-type classification loss, posterior modes arise naturally.

%%%%%%%%%%%%%%%%%%%%%%%%%%%%%%%%%%%%%%%%%%%%%%%%%%%%%%%%%%%%
\subsection{Admissibility}
\label{subsec:admissibility}
%%%%%%%%%%%%%%%%%%%%%%%%%%%%%%%%%%%%%%%%%%%%%%%%%%%%%%%%%%%%

A decision rule is inadmissible if another rule has risk no greater for
every \(\theta\) and strictly smaller for at least one \(\theta\).

A rule is admissible if no such uniformly better competitor exists.

Admissibility is therefore a non-dominance criterion.

%%%%%%%%%%%%%%%%%%%%%%%%%%%%%%%%%%%%%%%%%%%%%%%%%%%%%%%%%%%%
\subsection{Minimax Rules}
\label{subsec:minimax-rules}
%%%%%%%%%%%%%%%%%%%%%%%%%%%%%%%%%%%%%%%%%%%%%%%%%%%%%%%%%%%%

A minimax rule minimizes the worst-case risk:

\[
\boxed{
\delta^*
=
\arg\min_\delta
\sup_{\theta\in\Theta}
R
(
\theta,\delta
).
}
\]

Minimax procedures protect against the most unfavorable parameter value.

They are especially useful when no prior distribution is available.

%%%%%%%%%%%%%%%%%%%%%%%%%%%%%%%%%%%%%%%%%%%%%%%%%%%%%%%%%%%%
\subsection{Bayes Versus Minimax}
\label{subsec:bayes-versus-minimax}
%%%%%%%%%%%%%%%%%%%%%%%%%%%%%%%%%%%%%%%%%%%%%%%%%%%%%%%%%%%%

Bayes procedures minimize average risk under a prior.

Minimax procedures minimize maximum risk.

Thus:

\[
\boxed{
\text{Bayes}
=
\text{average-case under a prior},
}
\]

while

\[
\boxed{
\text{minimax}
=
\text{worst-case protection}.
}
\]

Some procedures can be both Bayes and minimax.

%%%%%%%%%%%%%%%%%%%%%%%%%%%%%%%%%%%%%%%%%%%%%%%%%%%%%%%%%%%%
\subsection{James--Stein Phenomenon Preview}
\label{subsec:james-stein-preview}
%%%%%%%%%%%%%%%%%%%%%%%%%%%%%%%%%%%%%%%%%%%%%%%%%%%%%%%%%%%%

In dimension

\[
p\geq3,
\]

estimating a multivariate normal mean under total squared-error loss
reveals a surprising phenomenon.

The ordinary coordinatewise sample mean can be inadmissible.

Shrinkage estimators can achieve lower total risk.

This result motivates modern regularization and shrinkage methods.

%%%%%%%%%%%%%%%%%%%%%%%%%%%%%%%%%%%%%%%%%%%%%%%%%%%%%%%%%%%%
\subsection{Exact Versus Asymptotic Inference}
\label{subsec:exact-versus-asymptotic-mathstat}
%%%%%%%%%%%%%%%%%%%%%%%%%%%%%%%%%%%%%%%%%%%%%%%%%%%%%%%%%%%%

Exact inference uses a finite-sample distribution known under the model.

Examples include:

\[
\boxed{
t,
}
\]

\[
\boxed{
\chi^2,
}
\]

and

\[
\boxed{
F
}
\]

results under normal sampling.

Asymptotic inference uses large-sample approximations.

Examples include:

\[
\boxed{
\text{Wald},
}
\]

\[
\boxed{
\text{score},
}
\]

and

\[
\boxed{
\text{likelihood-ratio}
}
\]

procedures in general parametric models.

%%%%%%%%%%%%%%%%%%%%%%%%%%%%%%%%%%%%%%%%%%%%%%%%%%%%%%%%%%%%
\subsection{Regularity Conditions}
\label{subsec:regularity-conditions}
%%%%%%%%%%%%%%%%%%%%%%%%%%%%%%%%%%%%%%%%%%%%%%%%%%%%%%%%%%%%

Many asymptotic results rely on assumptions such as:

\[
\boxed{
\text{identifiability},
}
\]

\[
\boxed{
\text{smooth likelihoods},
}
\]

\[
\boxed{
\text{interior parameter values},
}
\]

\[
\boxed{
\text{finite information},
}
\]

and

\[
\boxed{
\text{valid interchange of differentiation and integration}.
}
\]

When these fail, classical asymptotic results can fail.

%%%%%%%%%%%%%%%%%%%%%%%%%%%%%%%%%%%%%%%%%%%%%%%%%%%%%%%%%%%%
\subsection{Boundary Problems}
\label{subsec:parameter-boundary-problems}
%%%%%%%%%%%%%%%%%%%%%%%%%%%%%%%%%%%%%%%%%%%%%%%%%%%%%%%%%%%%

Suppose a null parameter lies on the boundary of the parameter space.

For example,

\[
\sigma^2=0
\]

may lie on the boundary of

\[
\sigma^2\geq0.
\]

Then Wilks' chi-square result may not hold in its ordinary form.

Boundary problems often require specialized reference distributions.

%%%%%%%%%%%%%%%%%%%%%%%%%%%%%%%%%%%%%%%%%%%%%%%%%%%%%%%%%%%%
\subsection{Nuisance Parameters}
\label{subsec:nuisance-parameters-mathstat}
%%%%%%%%%%%%%%%%%%%%%%%%%%%%%%%%%%%%%%%%%%%%%%%%%%%%%%%%%%%%

A nuisance parameter is an unknown parameter required by the model but
not the primary scientific target.

For example, when estimating normal mean \(\mu\), the variance
\(\sigma^2\) may be a nuisance parameter.

Statistical procedures may handle nuisance parameters through:

\[
\boxed{
\text{profiling},
}
\]

\[
\boxed{
\text{conditioning},
}
\]

\[
\boxed{
\text{marginalization},
}
\]

or

\[
\boxed{
\text{plug-in estimation}.
}
\]

%%%%%%%%%%%%%%%%%%%%%%%%%%%%%%%%%%%%%%%%%%%%%%%%%%%%%%%%%%%%
\subsection{Profile Likelihood}
\label{subsec:profile-likelihood}
%%%%%%%%%%%%%%%%%%%%%%%%%%%%%%%%%%%%%%%%%%%%%%%%%%%%%%%%%%%%

Suppose the parameter is

\[
(\psi,\lambda),
\]

where \(\psi\) is the parameter of interest and \(\lambda\) is nuisance.

For each fixed \(\psi\), maximize over \(\lambda\):

\[
\boxed{
L_p(\psi)
=
\sup_\lambda
L(\psi,\lambda).
}
\]

This is the profile likelihood.

It reduces a multidimensional likelihood to a function of the parameter
of interest.

%%%%%%%%%%%%%%%%%%%%%%%%%%%%%%%%%%%%%%%%%%%%%%%%%%%%%%%%%%%%
\subsection{Likelihood-Based Confidence Intervals}
\label{subsec:likelihood-based-confidence-intervals}
%%%%%%%%%%%%%%%%%%%%%%%%%%%%%%%%%%%%%%%%%%%%%%%%%%%%%%%%%%%%

An approximate likelihood-ratio confidence set for scalar \(\theta\) may
be written

\[
\boxed{
2
[
\ell(\widehat{\theta})
-
\ell(\theta)
]
\leq
\chi^2_{1,1-\alpha}.
}
\]

Unlike a Wald interval, this interval need not be symmetric around the
MLE.

Likelihood-based intervals can therefore behave better for skewed
parameter problems.

%%%%%%%%%%%%%%%%%%%%%%%%%%%%%%%%%%%%%%%%%%%%%%%%%%%%%%%%%%%%
\subsection{Information Matrix for Vector Parameters}
\label{subsec:information-matrix}
%%%%%%%%%%%%%%%%%%%%%%%%%%%%%%%%%%%%%%%%%%%%%%%%%%%%%%%%%%%%

For parameter vector

\[
\boldsymbol{\theta}\in\mathbb{R}^k,
\]

the score vector is

\[
U(\boldsymbol{\theta})
=
\nabla
\ell(\boldsymbol{\theta}).
\]

The Fisher information matrix is

\[
\boxed{
I(\boldsymbol{\theta})
=
\mathbb{E}
[
U(\boldsymbol{\theta})
U(\boldsymbol{\theta})^T
].
}
\]

Under regularity conditions,

\[
\boxed{
I(\boldsymbol{\theta})
=
-
\mathbb{E}
[
\nabla^2
\ell(\boldsymbol{\theta})
].
}
\]

%%%%%%%%%%%%%%%%%%%%%%%%%%%%%%%%%%%%%%%%%%%%%%%%%%%%%%%%%%%%
\subsection{Multivariate Cramér--Rao Bound}
\label{subsec:multivariate-cramer-rao}
%%%%%%%%%%%%%%%%%%%%%%%%%%%%%%%%%%%%%%%%%%%%%%%%%%%%%%%%%%%%

For an unbiased estimator vector under suitable regularity conditions,

\[
\boxed{
\operatorname{Cov}
(
\widehat{\boldsymbol{\theta}}
)
-
I(\boldsymbol{\theta})^{-1}
}
\]

is positive semidefinite.

Thus

\[
\boxed{
I(\boldsymbol{\theta})^{-1}
}
\]

acts as a matrix lower bound on covariance.

%%%%%%%%%%%%%%%%%%%%%%%%%%%%%%%%%%%%%%%%%%%%%%%%%%%%%%%%%%%%
\subsection{Estimating Equations Preview}
\label{subsec:estimating-equations-preview}
%%%%%%%%%%%%%%%%%%%%%%%%%%%%%%%%%%%%%%%%%%%%%%%%%%%%%%%%%%%%

Not every estimator arises by maximizing a likelihood.

An estimating equation has the form

\[
\boxed{
\sum_{i=1}^{n}
\psi
(
X_i,\theta
)
=
0.
}
\]

Solving this equation yields an estimator.

Examples include:

\[
\boxed{
\text{method of moments},
}
\]

\[
\boxed{
\text{score equations},
}
\]

and many robust estimators.

%%%%%%%%%%%%%%%%%%%%%%%%%%%%%%%%%%%%%%%%%%%%%%%%%%%%%%%%%%%%
\subsection{M-Estimators Preview}
\label{subsec:m-estimators-preview}
%%%%%%%%%%%%%%%%%%%%%%%%%%%%%%%%%%%%%%%%%%%%%%%%%%%%%%%%%%%%

An M-estimator minimizes an objective

\[
\boxed{
\sum_{i=1}^{n}
\rho
(
X_i,\theta
)
}
\]

or equivalently solves an estimating equation based on

\[
\psi
=
\frac{
\partial\rho
}{
\partial\theta
}.
\]

Maximum likelihood and least squares are special cases.

Robust statistics often replaces squared error with less outlier-sensitive
loss functions.

%%%%%%%%%%%%%%%%%%%%%%%%%%%%%%%%%%%%%%%%%%%%%%%%%%%%%%%%%%%%
\subsection{Influence Functions Preview}
\label{subsec:influence-functions-preview}
%%%%%%%%%%%%%%%%%%%%%%%%%%%%%%%%%%%%%%%%%%%%%%%%%%%%%%%%%%%%

An influence function measures the infinitesimal effect of contamination
at a point on a statistical functional.

Roughly,

\[
\boxed{
IF(x;T,F)
=
\lim_{\varepsilon\to0}
\frac{
T
(
(1-\varepsilon)F+\varepsilon\Delta_x
)
-
T(F)
}{
\varepsilon
}.
}
\]

Bounded influence is one indication of robustness.

%%%%%%%%%%%%%%%%%%%%%%%%%%%%%%%%%%%%%%%%%%%%%%%%%%%%%%%%%%%%
\subsection{Empirical Process Preview}
\label{subsec:empirical-process-preview-mathstat}
%%%%%%%%%%%%%%%%%%%%%%%%%%%%%%%%%%%%%%%%%%%%%%%%%%%%%%%%%%%%

The empirical distribution function is

\[
F_n(x).
\]

The centered and scaled empirical process is

\[
\boxed{
\sqrt{n}
[
F_n(x)-F(x)
].
}
\]

Its asymptotic behavior underlies:

\[
\boxed{
\text{Kolmogorov--Smirnov theory},
}
\]

\[
\boxed{
\text{nonparametric confidence bands},
}
\]

and many modern asymptotic results.

%%%%%%%%%%%%%%%%%%%%%%%%%%%%%%%%%%%%%%%%%%%%%%%%%%%%%%%%%%%%
\subsection{Statistical Functional}
\label{subsec:statistical-functional}
%%%%%%%%%%%%%%%%%%%%%%%%%%%%%%%%%%%%%%%%%%%%%%%%%%%%%%%%%%%%

Many estimands can be viewed as functions of a population distribution
\(F\).

For example,

\[
\boxed{
T(F)
=
\int x\,dF(x)
}
\]

is the mean functional.

The empirical estimator replaces \(F\) with \(F_n\):

\[
\boxed{
T(F_n).
}
\]

This viewpoint unifies parametric, nonparametric, and bootstrap methods.

%%%%%%%%%%%%%%%%%%%%%%%%%%%%%%%%%%%%%%%%%%%%%%%%%%%%%%%%%%%%
\subsection{Bootstrap Through Mathematical Statistics}
\label{subsec:bootstrap-mathstat}
%%%%%%%%%%%%%%%%%%%%%%%%%%%%%%%%%%%%%%%%%%%%%%%%%%%%%%%%%%%%

The nonparametric bootstrap treats the empirical distribution

\[
F_n
\]

as an estimate of \(F\).

A bootstrap sample is drawn from \(F_n\).

The distribution of

\[
\boxed{
T(F_n^*)
}
\]

conditional on the observed data is used to approximate the sampling
distribution of

\[
T(F_n).
\]

Computational Statistics develops this more fully.

%%%%%%%%%%%%%%%%%%%%%%%%%%%%%%%%%%%%%%%%%%%%%%%%%%%%%%%%%%%%
\subsection{Mathematical Statistics Workflow}
\label{subsec:mathematical-statistics-workflow}
%%%%%%%%%%%%%%%%%%%%%%%%%%%%%%%%%%%%%%%%%%%%%%%%%%%%%%%%%%%%

A theoretical statistical analysis often follows the sequence:

\begin{enumerate}
    \item specify the sample space;
    \item specify the statistical model;
    \item identify the parameter space;
    \item verify identifiability;
    \item define the estimand;
    \item choose an estimator or decision rule;
    \item derive its sampling distribution when possible;
    \item compute bias, variance, or risk;
    \item identify sufficient statistics;
    \item derive likelihood and score functions;
    \item calculate Fisher information;
    \item determine exact or asymptotic distributions;
    \item construct tests or confidence sets;
    \item analyze large-sample properties;
    \item evaluate optimality or robustness.
\end{enumerate}

%%%%%%%%%%%%%%%%%%%%%%%%%%%%%%%%%%%%%%%%%%%%%%%%%%%%%%%%%%%%
\subsection{Common Errors}
\label{subsec:mathematical-statistics-common-errors}
%%%%%%%%%%%%%%%%%%%%%%%%%%%%%%%%%%%%%%%%%%%%%%%%%%%%%%%%%%%%

\subsubsection{Confusing a Parameter with an Estimator}

The parameter is fixed within the frequentist model.

The estimator is random because it depends on random data.

%%%%%%%%%%%%%%%%%%%%%%%%%%%%%%%%%%%%%%%%%%%%%%%%%%%%%%%%%%%%
\subsubsection{Assuming Unbiased Means Best}

Unbiasedness is only one property.

A biased estimator may have substantially lower mean squared error.

%%%%%%%%%%%%%%%%%%%%%%%%%%%%%%%%%%%%%%%%%%%%%%%%%%%%%%%%%%%%
\subsubsection{Assuming Consistency Means Unbiasedness}

An estimator can be biased for every finite sample size and still be
consistent.

Consistency concerns behavior as

\[
n\to\infty.
\]

%%%%%%%%%%%%%%%%%%%%%%%%%%%%%%%%%%%%%%%%%%%%%%%%%%%%%%%%%%%%
\subsubsection{Assuming an MLE Is Always Unbiased}

Maximum likelihood is defined by likelihood maximization, not
unbiasedness.

%%%%%%%%%%%%%%%%%%%%%%%%%%%%%%%%%%%%%%%%%%%%%%%%%%%%%%%%%%%%
\subsubsection{Solving Only the Score Equation Without Checking the Maximum}

The equation

\[
U(\theta)=0
\]

identifies stationary points.

One must still consider:

\[
\boxed{
\text{boundaries},
}
\]

\[
\boxed{
\text{second derivatives},
}
\]

or

\[
\boxed{
\text{global likelihood behavior}.
}
\]

%%%%%%%%%%%%%%%%%%%%%%%%%%%%%%%%%%%%%%%%%%%%%%%%%%%%%%%%%%%%
\subsubsection{Ignoring the Parameter Space}

A formal solution may lie outside the allowed parameter space.

Maximum likelihood optimization must respect

\[
\theta\in\Theta.
\]

%%%%%%%%%%%%%%%%%%%%%%%%%%%%%%%%%%%%%%%%%%%%%%%%%%%%%%%%%%%%
\subsubsection{Applying Cramér--Rao Without Checking Conditions}

The ordinary Cramér--Rao bound requires regularity assumptions and
unbiasedness in its standard form.

It does not apply mechanically to every estimation problem.

%%%%%%%%%%%%%%%%%%%%%%%%%%%%%%%%%%%%%%%%%%%%%%%%%%%%%%%%%%%%
\subsubsection{Assuming Sufficiency Means the Statistic Is a Good Estimator}

Sufficiency concerns information preservation.

A sufficient statistic is not necessarily itself an estimator of the
target parameter.

%%%%%%%%%%%%%%%%%%%%%%%%%%%%%%%%%%%%%%%%%%%%%%%%%%%%%%%%%%%%
\subsubsection{Confusing Sufficiency and Completeness}

Sufficiency concerns retaining parameter information.

Completeness concerns uniqueness of zero-mean functions of a statistic.

These are different properties.

%%%%%%%%%%%%%%%%%%%%%%%%%%%%%%%%%%%%%%%%%%%%%%%%%%%%%%%%%%%%
\subsubsection{Assuming Fisher Information Is Observed Curvature}

Expected Fisher information and observed information are related but not
identical.

One is an expectation under the model; the other depends on the observed
sample.

%%%%%%%%%%%%%%%%%%%%%%%%%%%%%%%%%%%%%%%%%%%%%%%%%%%%%%%%%%%%
\subsubsection{Using Large-Sample Theory in Very Small Samples Without Caution}

Asymptotic approximations may be inaccurate when the sample is small,
parameters are near boundaries, or likelihoods are highly skewed.

%%%%%%%%%%%%%%%%%%%%%%%%%%%%%%%%%%%%%%%%%%%%%%%%%%%%%%%%%%%%
\subsubsection{Assuming Wilks' Theorem Always Gives Chi-Square}

Nonidentifiability, boundary parameters, mixture models, and other
irregular situations can invalidate the usual chi-square limit.

%%%%%%%%%%%%%%%%%%%%%%%%%%%%%%%%%%%%%%%%%%%%%%%%%%%%%%%%%%%%
\subsubsection{Interpreting Failure to Reject as Proof of the Null}

Hypothesis tests control error rates under repeated sampling.

They do not assign posterior probabilities to hypotheses.

%%%%%%%%%%%%%%%%%%%%%%%%%%%%%%%%%%%%%%%%%%%%%%%%%%%%%%%%%%%%
\subsubsection{Ignoring Nuisance Parameters}

Uncertainty in nuisance parameters can materially affect inference for
the parameter of interest.

%%%%%%%%%%%%%%%%%%%%%%%%%%%%%%%%%%%%%%%%%%%%%%%%%%%%%%%%%%%%
\subsubsection{Comparing Estimators Only by Variance When They Are Biased}

For biased estimators, mean squared error is often the more appropriate
criterion.

%%%%%%%%%%%%%%%%%%%%%%%%%%%%%%%%%%%%%%%%%%%%%%%%%%%%%%%%%%%%
\subsubsection{Ignoring the Loss Function in Decision Problems}

A decision rule can only be called optimal relative to a specified loss or
risk criterion.

Different losses can produce different optimal rules.

%%%%%%%%%%%%%%%%%%%%%%%%%%%%%%%%%%%%%%%%%%%%%%%%%%%%%%%%%%%%
\subsection{Applications}
\label{subsec:mathematical-statistics-applications}
%%%%%%%%%%%%%%%%%%%%%%%%%%%%%%%%%%%%%%%%%%%%%%%%%%%%%%%%%%%%

\begin{applicationbox}

\textbf{Scientific Inference.}
Mathematical statistics supplies rigorous foundations for parameter
estimation, uncertainty quantification, and hypothesis testing.

\medskip

\textbf{Engineering.}
Likelihood, information, and optimal estimation theory are used in
reliability, signal processing, calibration, and quality control.

\medskip

\textbf{Biostatistics.}
Likelihood and asymptotic theory support generalized linear models,
survival models, clinical trials, and longitudinal analysis.

\medskip

\textbf{Econometrics.}
Consistency, asymptotic normality, estimating equations, and efficiency
are central to modern econometric inference.

\medskip

\textbf{Machine Learning.}
Risk minimization, likelihood, regularization, decision theory, and
asymptotics provide theoretical foundations for statistical learning.

\medskip

\textbf{Bayesian Statistics.}
Decision theory connects frequentist risk with posterior expected loss,
while likelihood provides the data component of Bayesian updating.

\medskip

\textbf{Computational Statistics.}
Many estimators cannot be solved analytically and require numerical
optimization, simulation, bootstrap, or Monte Carlo procedures.

\end{applicationbox}

%%%%%%%%%%%%%%%%%%%%%%%%%%%%%%%%%%%%%%%%%%%%%%%%%%%%%%%%%%%%
\subsection{Exercises}
\label{subsec:mathematical-statistics-exercises}
%%%%%%%%%%%%%%%%%%%%%%%%%%%%%%%%%%%%%%%%%%%%%%%%%%%%%%%%%%%%

\begin{exercise}[1]
Define a statistical model.
\end{exercise}

\begin{exercise}[1]
Define a parameter space.
\end{exercise}

\begin{exercise}[1]
Distinguish a parameter from a statistic.
\end{exercise}

\begin{exercise}[1]
Distinguish an estimator from an estimate.
\end{exercise}

\begin{exercise}[1]
Define identifiability.
\end{exercise}

\begin{exercise}[1]
Define estimator bias.
\end{exercise}

\begin{exercise}[1]
Define mean squared error.
\end{exercise}

\begin{exercise}[1]
Define consistency.
\end{exercise}

\begin{exercise}[1]
Define a likelihood function.
\end{exercise}

\begin{exercise}[1]
Define a maximum likelihood estimator.
\end{exercise}

\begin{exercise}[1]
Define a score function.
\end{exercise}

\begin{exercise}[1]
Define Fisher information.
\end{exercise}

\begin{exercise}[1]
Define a sufficient statistic.
\end{exercise}

\begin{exercise}[1]
Define a complete statistic.
\end{exercise}

\begin{exercise}[1]
Define a pivotal quantity.
\end{exercise}

\begin{exercise}[1]
Define a statistical loss function.
\end{exercise}

\begin{exercise}[2]
Show that

\[
MSE(\widehat{\theta})
=
\operatorname{Var}(\widehat{\theta})
+
\operatorname{Bias}(\widehat{\theta})^2.
\]
\end{exercise}

\begin{exercise}[2]
Suppose

\[
\mathbb{E}[\widehat{\theta}]
=
\theta+\frac1n.
\]

Find the bias.
\end{exercise}

\begin{exercise}[2]
Suppose an estimator has variance

\[
\frac4n
\]

and bias

\[
\frac1n.
\]

Find its MSE.
\end{exercise}

\begin{exercise}[2]
For Bernoulli data, derive the method-of-moments estimator of \(p\).
\end{exercise}

\begin{exercise}[2]
For Poisson data, derive the method-of-moments estimator of \(\lambda\).
\end{exercise}

\begin{exercise}[2]
For exponential data, derive the MLE of the rate parameter.
\end{exercise}

\begin{exercise}[2]
For Bernoulli data, derive the MLE of \(p\).
\end{exercise}

\begin{exercise}[2]
For normal data with known variance, derive the MLE of \(\mu\).
\end{exercise}

\begin{exercise}[2]
For a Bernoulli observation, compute the score function.
\end{exercise}

\begin{exercise}[2]
For a Bernoulli observation, compute Fisher information.
\end{exercise}

\begin{exercise}[2]
For

\[
X_i\sim N(\mu,\sigma^2)
\]

with known \(\sigma^2\), compute Fisher information for \(\mu\).
\end{exercise}

\begin{exercise}[2]
Use the Cramér--Rao bound to show that the sample mean is efficient for
the normal mean with known variance.
\end{exercise}

\begin{exercise}[2]
Use the factorization theorem to show that

\[
\sum X_i
\]

is sufficient for a Poisson rate parameter.
\end{exercise}

\begin{exercise}[2]
State the Rao--Blackwell theorem.
\end{exercise}

\begin{exercise}[2]
State the Lehmann--Scheffé theorem.
\end{exercise}

\begin{exercise}[3]
Prove consistency of the sample mean using Chebyshev's inequality under
finite variance.
\end{exercise}

\begin{exercise}[3]
Show that convergence in mean square implies convergence in probability.
\end{exercise}

\begin{exercise}[3]
Derive the method-of-moments estimators for the parameters of a
two-parameter distribution of your choice.
\end{exercise}

\begin{exercise}[3]
Derive the MLEs of \(\mu\) and \(\sigma^2\) for a normal sample.
\end{exercise}

\begin{exercise}[3]
Show explicitly why the normal variance MLE is biased.
\end{exercise}

\begin{exercise}[3]
Prove the invariance property of maximum likelihood.
\end{exercise}

\begin{exercise}[3]
Prove that the expected score is zero under regularity conditions.
\end{exercise}

\begin{exercise}[3]
Show the equivalence

\[
I(\theta)
=
\mathbb{E}[U(\theta)^2]
=
-
\mathbb{E}[\ell''(\theta)]
\]

under appropriate regularity conditions.
\end{exercise}

\begin{exercise}[3]
Derive the Cramér--Rao lower bound.
\end{exercise}

\begin{exercise}[3]
Use the factorization theorem to identify a sufficient statistic for the
exponential rate parameter.
\end{exercise}

\begin{exercise}[3]
Show that

\[
(
\overline{X},S^2
)
\]

is sufficient for \((\mu,\sigma^2)\) in the normal family.
\end{exercise}

\begin{exercise}[3]
Prove the Rao--Blackwell theorem using the law of total variance.
\end{exercise}

\begin{exercise}[3]
Explain why completeness leads to uniqueness of unbiased estimators based
on a complete sufficient statistic.
\end{exercise}

\begin{exercise}[3]
Rewrite a Bernoulli distribution in exponential-family form.
\end{exercise}

\begin{exercise}[3]
Rewrite a Poisson distribution in exponential-family form.
\end{exercise}

\begin{exercise}[3]
Identify the natural parameter and sufficient statistic in the Poisson
family.
\end{exercise}

\begin{exercise}[3]
Show that derivatives of the log-partition function generate moments of
the natural statistic.
\end{exercise}

\begin{exercise}[3]
Construct a confidence interval by inverting a normal pivot.
\end{exercise}

\begin{exercise}[3]
Construct an exact variance confidence interval using the chi-square
pivot.
\end{exercise}

\begin{exercise}[3]
Explain the duality between tests and confidence sets.
\end{exercise}

\begin{exercise}[3]
Derive a likelihood-ratio test for a Bernoulli parameter.
\end{exercise}

\begin{exercise}[3]
Compare Wald, score, and likelihood-ratio tests for a one-parameter
model.
\end{exercise}

\begin{exercise}[3]
Apply the delta method to

\[
g(\theta)=\theta^2.
\]
\end{exercise}

\begin{exercise}[3]
Apply the delta method to

\[
g(\theta)=1/\theta.
\]
\end{exercise}

\begin{exercise}[3]
Use Slutsky's theorem to justify replacing a population standard
deviation by a consistent sample estimate in an asymptotic statistic.
\end{exercise}

\begin{exercise}[3]
Derive the distribution of the sample maximum from a continuous CDF.
\end{exercise}

\begin{exercise}[3]
Derive the density of the \(k\)-th order statistic.
\end{exercise}

\begin{exercise}[3]
Show that

\[
F(X_{(k)})
\]

has a beta distribution under a continuous model.
\end{exercise}

\begin{exercise}[3]
Show that squared-error risk equals MSE for point estimation.
\end{exercise}

\begin{exercise}[3]
Prove that the posterior mean minimizes posterior expected squared-error
loss.
\end{exercise}

\begin{exercise}[3]
Explain the difference between Bayes and minimax decision rules.
\end{exercise}

\begin{exercise}[4]
Prove the Neyman--Pearson lemma.
\end{exercise}

\begin{exercise}[4]
Prove the Karlin--Rubin theorem.
\end{exercise}

\begin{exercise}[4]
Study complete sufficient statistics in one-parameter exponential
families.
\end{exercise}

\begin{exercise}[4]
Derive a UMVU estimator using the Lehmann--Scheffé theorem.
\end{exercise}

\begin{exercise}[4]
Study Basu's theorem relating complete sufficient and ancillary
statistics.
\end{exercise}

\begin{exercise}[4]
Derive the asymptotic normality of the MLE.
\end{exercise}

\begin{exercise}[4]
Derive Wilks' theorem under regularity conditions.
\end{exercise}

\begin{exercise}[4]
Study local alternatives and asymptotic power of Wald, score, and
likelihood-ratio tests.
\end{exercise}

\begin{exercise}[4]
Derive the multivariate Cramér--Rao bound.
\end{exercise}

\begin{exercise}[4]
Study profile likelihood inference with nuisance parameters.
\end{exercise}

\begin{exercise}[4]
Investigate irregular likelihood problems at parameter-space boundaries.
\end{exercise}

\begin{exercise}[4]
Study the asymptotic efficiency of method-of-moments and maximum
likelihood estimators.
\end{exercise}

\begin{exercise}[4]
Derive the multivariate delta method.
\end{exercise}

\begin{exercise}[4]
Study M-estimators and their asymptotic distributions.
\end{exercise}

\begin{exercise}[4]
Investigate influence functions and robust asymptotic theory.
\end{exercise}

\begin{exercise}[4]
Study the James--Stein estimator and prove domination of the sample mean
under appropriate squared-error loss.
\end{exercise}

\begin{exercise}[4]
Study admissibility and complete-class theorems.
\end{exercise}

\begin{exercise}[4]
Investigate minimax estimation using least-favorable priors.
\end{exercise}

\begin{exercise}[4]
Study empirical-process convergence and Brownian bridges.
\end{exercise}

\begin{exercise}[4]
Investigate theoretical consistency of the nonparametric bootstrap.
\end{exercise}

%%%%%%%%%%%%%%%%%%%%%%%%%%%%%%%%%%%%%%%%%%%%%%%%%%%%%%%%%%%%
\subsection{Conceptual Summary}
\label{subsec:mathematical-statistics-summary}
%%%%%%%%%%%%%%%%%%%%%%%%%%%%%%%%%%%%%%%%%%%%%%%%%%%%%%%%%%%%

A statistical model is

\[
\boxed{
\mathcal{P}
=
\{
P_\theta:
\theta\in\Theta
\}.
}
\]

An estimator

\[
\widehat{\theta}
\]

is evaluated through properties such as

\[
\boxed{
\text{bias},
}
\]

\[
\boxed{
\text{variance},
}
\]

\[
\boxed{
MSE,
}
\]

\[
\boxed{
\text{consistency},
}
\]

and

\[
\boxed{
\text{efficiency}.
}
\]

Mean squared error satisfies

\[
\boxed{
MSE
=
\operatorname{Var}
+
\operatorname{Bias}^2.
}
\]

Maximum likelihood uses

\[
\boxed{
\widehat{\theta}_{MLE}
=
\arg\max_{\theta\in\Theta}
L(\theta).
}
\]

The score is

\[
\boxed{
U(\theta)
=
\ell'(\theta),
}
\]

and Fisher information is

\[
\boxed{
I(\theta)
=
\mathbb{E}
[
U(\theta)^2
].
}
\]

The Cramér--Rao inequality gives

\[
\boxed{
\operatorname{Var}
(
\widehat{\theta}
)
\geq
\frac1{
I_n(\theta)
}
}
\]

for regular unbiased estimation.

Sufficiency can be identified using

\[
\boxed{
f_\theta(x)
=
g_\theta(T(x))h(x).
}
\]

The Rao--Blackwell and Lehmann--Scheffé theorems connect sufficiency,
completeness, and optimal unbiased estimation.

Likelihood testing uses

\[
\boxed{
\Lambda
=
\frac{
\sup_{\Theta_0}L
}{
\sup_{\Theta}L
}.
}
\]

Under regularity conditions,

\[
\boxed{
-2\log\Lambda
\xrightarrow{d}
\chi^2.
}
\]

The MLE commonly satisfies

\[
\boxed{
\sqrt{n}
(
\widehat{\theta}-\theta
)
\xrightarrow{d}
N
\left(
0,
I_1(\theta)^{-1}
\right).
}
\]

Mathematical statistics therefore unifies

\[
\boxed{
\text{probability}
+
\text{estimation}
+
\text{likelihood}
+
\text{testing}
+
\text{asymptotics}
+
\text{decision theory}.
}
\]

%%%%%%%%%%%%%%%%%%%%%%%%%%%%%%%%%%%%%%%%%%%%%%%%%%%%%%%%%%%%
\subsection{Section Connections}
\label{subsec:mathematical-statistics-connections}
%%%%%%%%%%%%%%%%%%%%%%%%%%%%%%%%%%%%%%%%%%%%%%%%%%%%%%%%%%%%

\chapterconnection{
Mathematical Statistics provides the theoretical foundation for the
estimation, confidence interval, hypothesis testing, regression, ANOVA,
categorical-data, nonparametric, and multivariate procedures developed
throughout Chapter~\ref{chap:statistics}. Sufficiency and information
explain how data encode parameter information, likelihood theory unifies
many estimators and tests, and asymptotic theory explains why normal and
chi-square approximations repeatedly appear in statistical practice.
Decision theory introduces loss, risk, Bayes rules, admissibility, and
minimaxity. The next section, Bayesian Statistics, develops probability
models in which unknown parameters themselves are assigned prior
distributions and inference proceeds through posterior distributions.
}

%%%%%%%%%%%%%%%%%%%%%%%%%%%%%%%%%%%%%%%%%%%%%%%%%%%%%%%%%%%%
% SECTION SOURCE TRACKING
%%%%%%%%%%%%%%%%%%%%%%%%%%%%%%%%%%%%%%%%%%%%%%%%%%%%%%%%%%%%

%------------------------------------------------------------
% 8.13 Mathematical Statistics
%------------------------------------------------------------
%
% Source basis:
%   - Standard mathematical statistics.
%   - Standard likelihood and estimation theory.
%   - Standard frequentist decision theory and asymptotic inference.
%
% Used for:
%   - statistical models and parameter spaces;
%   - identifiability;
%   - estimators and sampling distributions;
%   - bias;
%   - variance;
%   - mean squared error;
%   - consistency;
%   - asymptotic normality;
%   - method of moments;
%   - maximum likelihood;
%   - score functions;
%   - Fisher information;
%   - Cramer--Rao lower bounds;
%   - efficiency;
%   - sufficient statistics;
%   - factorization theorem;
%   - minimal sufficiency;
%   - ancillary statistics;
%   - completeness;
%   - Rao--Blackwell theorem;
%   - Lehmann--Scheffe theorem;
%   - UMVU estimation;
%   - exponential families;
%   - pivots;
%   - confidence intervals by pivoting;
%   - confidence sets by test inversion;
%   - Neyman--Pearson theory;
%   - most powerful and UMP tests;
%   - monotone likelihood ratio;
%   - likelihood-ratio tests;
%   - Wald and score tests;
%   - Wilks theorem;
%   - asymptotic MLE theory;
%   - observed information;
%   - delta method;
%   - multivariate delta method;
%   - Slutsky theorem;
%   - continuous mapping theorem;
%   - asymptotic efficiency;
%   - order statistics;
%   - decision theory;
%   - loss and risk;
%   - Bayes risk;
%   - admissibility;
%   - minimaxity;
%   - James--Stein preview;
%   - nuisance parameters;
%   - profile likelihood;
%   - likelihood confidence intervals;
%   - information matrices;
%   - multivariate Cramer--Rao bounds;
%   - estimating equations;
%   - M-estimator previews;
%   - influence-function previews;
%   - empirical-process preview;
%   - bootstrap theory preview;
%   - applications and exercises.
%
% Notes:
%   - Bayesian Statistics is developed next in Section 8.14.
%   - Statistical Learning later develops risk minimization,
%     regularization, and high-dimensional estimation more deeply.
%   - Computational Statistics develops numerical likelihood,
%     bootstrap, simulation, and Monte Carlo methods in greater detail.
%
%%%%%%%%%%%%%%%%%%%%%%%%%%%%%%%%%%%%%%%%%%%%%%%%%%%%%%%%%%%%